% ContEst.tex -- Control and Estimation Co-Design
% Automatica style (autart.cls, two-column). Compile with the autart
% class from the Automatica sample. Per Automatica policy, NO author
% macros are defined in this file.
%
% Bibliography keys come from:
%   (i)  your References.bib:
%        garcia2019control, feldbaum1960dual, bar1974dual,
%        kumar2015stochastic, resnick2019probability, anderson2012optimal
%   (ii) ContEst_related_work.bib (generated alongside this file)
%   (iii) two works you own; suggested BibTeX in the comment block at the
%         end:  grigoriadis1998integrate (paper 2),
%               ramadan_ekfkoopman       (paper 3, eKF--Koopman).
% Make sure every key resolves before running bibtex.

\documentclass[twocolumn,amsthm]{autart}


\usepackage[utf8]{inputenc}
\usepackage{comment}
\usepackage{algorithm}
\newtheorem{definition}{Definition}[section]
\usepackage{algpseudocode}
\newcommand{\Y}{\mathcal{Y}}
\newcommand{\Ll}{\mathcal{L}}
\newcommand{\E}{\mathbb{E\,}}
\newcommand{\Cov}{\textbf{Cov \,}}
\usepackage{graphicx}
\usepackage[version=4]{mhchem}
\usepackage{siunitx}
\usepackage{longtable,tabularx}
\usepackage{booktabs}
\usepackage{pifont}
\setlength\LTleft{0pt} 
\newcommand{\R}{\mathbb{R}}
\usepackage{graphics} % for pdf, bitmapped graphics files
\usepackage{float}
\usepackage{epsfig} % for postscript graphics files
%\usepackage{mathptmx} % assumes new font selection scheme installed
\usepackage{times} % assumes new font selection scheme installed
\usepackage{amsmath} % assumes amsmath package installed
\usepackage{amssymb}  % assumes amsmath package installed
\usepackage{mathtools}% http://ctan.org/pkg/mathtools
\usepackage{color}
\usepackage[dvipsnames]{xcolor}
\usepackage{hyperref}
% autart.cls does \@addtoreset{section}{part}, so newer hyperref builds section
% bookmark anchors as \theHpart.<n>. The class never defines \theHpart, so it was
% written unexpanded into output.out and broke the bookmark pass on Overleaf
% (Missing \endcsname / Use of \B_section.). Define it as a plain, fully
% \csname-expandable digit (this article never uses \part, so 0 is always right):
% this both writes correct anchors AND lets any already-cached bad .out read
% cleanly, so no "clear cached files" is required on Overleaf.
\providecommand*{\theHpart}{0}

\newtheorem{assumption}{Assumption}
\newtheorem{lemma}{Lemma}
\newtheorem{theorem}{Theorem}
\newtheorem{proposition}{Proposition}
\newtheorem{corollary}{Corollary}
%\hypersetup{
%    colorlinks=true,
%    linkcolor=blue,
%    filecolor=magenta,      
%    urlcolor=blue,
%}

\usepackage{bm}
\usepackage{times}
% Bob's additions
\usepackage{soul}
\newcommand{\cut}[1]{{\color{blue} \st{#1}}} % replace with {} to cut for real
\newcommand{\add}[1]{{\color{red} #1}} % replace blue with black to add for real
\newcommand{\concern}[1]{{\color{red} #1}} % replace blue with black to add for real
\newcommand{\matlab}{\mbox{\textsc{Matlab}}}

% ===================== Reviewer annotations (Mihai) =====================
% Color-coded todonotes: red = error/needs fixing, yellow = suggestion,
% blue = explanatory note (what the author did / how to read this).
% Remove this block and all \errornote/\suggestnote/\explainnote calls
% before final submission.
\usepackage[colorinlistoftodos,textsize=footnotesize]{todonotes}
% NOTE: autart is a two-column class with narrow margins, so notes are set
% "inline" (a colored box right in the text flow) rather than in the margin,
% which would get clipped. Switch to margin notes by removing "inline," if
% you compile single-column elsewhere.
\newcommand{\errornote}[1]{\todo[inline,color=red!40,bordercolor=red]{\textbf{ERROR:} #1}}
\newcommand{\suggestnote}[1]{\todo[inline,color=yellow!40,bordercolor=orange]{\textbf{Suggestion:} #1}}
\newcommand{\explainnote}[1]{\todo[inline,color=blue!25,bordercolor=blue]{\textbf{Note:} #1}}
% ===========================================================================

\begin{document}

\begin{frontmatter}

\title{Control and Estimation Co-Design via Envelope-Theorem Gradients\thanksref{footnoteinfo}}

 \thanks[footnoteinfo]{Corresponding author M.~S.~Ramadan. This material was
 based upon work supported by the U.S.\ Department of Energy, Office of Science,
 Office of Advanced Scientific Computing Research (ASCR), under Contract
 DE-AC02-06CH11347.}

\author[anl]{Mohammad S. Ramadan}\ead{mramadan@anl.gov},
\author[anl]{Philip Dinenis}\ead{pdinenis@anl.gov},
\author[anl]{Mihai Anitescu}\ead{anitescu@mcs.anl.gov}

\address[anl]{Mathematics and Computer Science Division, Argonne National
Laboratory, Lemont, IL 60439, USA}

\begin{keyword}
Control co-design; Stochastic optimal control; State estimation;
Semidefinite programming; Sensor placement; Sensitivity analysis.
\end{keyword}

\begin{abstract}
Instead of the sequential plant--control--estimator design pipeline, which is in general not optimal,
we propose the Control and Estimation Co-Design (ContEst) framework, which poses the entire system design as a single two-stage problem solved by first-order methods. A key
enabler is that the gradient of each inner control/estimation optimal cost with
respect to the design parameters is available from the dual variables
the inner solver already returns, by the envelope theorem: exactly under mild
regularity (a unique inner minimizer), and as a subgradient otherwise. That is, no differentiation
through the optimizer or its optimality conditions is needed. We first present the
framework in a linear quadratic Gaussian (LQG) regime, where the control and
estimation costs become two semidefinite programs (SDPs) whose design gradients are
read from the duals of the Lyapunov constraints. We extend the results to robust ($H_\infty$)
formulation to cover worst-case control and filtering problems. Furthermore, we present an extension to constrained and nonlinear systems, using an extended Kalman filter (eKF) and a local linearization of the information-state dynamics to yield an approximate but convex model predictive control (MPC) problem whose design gradients are read directly from the costates. We apply our methods to a variety of realistic design examples from different fields and report significant design improvements.
\end{abstract}

\end{frontmatter}

%=======================================================================
\section{Introduction}\label{sec:intro}
%=======================================================================
\explainnote{Roadmap of the whole paper, for a first-time reader (and for me while reviewing): Sec.~2 is related work; Sec.~3 sets up the general two-stage design problem; Sec.~4 specializes it to LQG, where the inner problems become the two SDPs~\eqref{eq:Jcont}--\eqref{eq:Jest}; Sec.~5 is the core theoretical contribution -- the envelope-theorem argument that makes the design gradient free, plus the regularity conditions (Assumption~\ref{as:env}) under which it's exact vs.\ merely a subgradient; Sec.~6 redoes the same story for worst-case $H_\infty$; Sec.~7 covers the eKF/MPC nonlinear extension; Sec.~8 is a complexity discussion (see the section-ordering note there); Sec.~9 is the four numerical case studies; the appendices carry the heavier derivations (information-state reduction, the LQG regularity proof, and the subgradient existence proof).}

Engineering systems are still, overwhelmingly, designed sequentially. A
plant is first conceived by some set of disciplines, the sensing and
instrumentation are specified next, and the controller and estimator are
designed last, once the plant, actuation and the measurement chain is frozen
\cite{garcia2019control}. Each stage restricts the next, so decisions made
early on, before any closed-loop consideration, can bound the performance
achievable at the end. This division of design effort is convenient, but
it is in general not optimal: the sequential solution (if exists) is one feasible
point of a larger joint problem and need not be an optimal or even a stationary point of it \cite{fathy2001coupling}.

Control co-design (CCD) addresses this on the control side, optimizing the
plant and the controller together~\cite{garcia2019control,allison2014multidisciplinary}. The same idea appeared in
process systems engineering as integrated design and control \cite{sakizlis2004recent,sharifzadeh2013integration} and in aerospace as
simultaneous structure/control design \cite{onoda1987approach,hale1985optimal}. What these literatures share is a
plant-controller scope: the sensing configuration and the estimator design are
typically fixed, or postponed to later steps. Yet for any
system without full and noiseless measurements (i.e. almost every real
system), achievable control performance is limited by the sensing and estimation quality.

This paper proposes the Control and Estimation Co-Design (ContEst) framework, in which the design parameters governing actuation and sensing/estimation are placed in a single two-stage optimization whose second-stage objective is the stochastic optimal control (dual control) cost. Because the system is
output feedback, the cost depends on the system not through the unobserved state but through its information state (commonly the filtering density
\cite{kumar2015stochastic}), whose evolution couples the output
equation and the state dynamics; optimizing through it therefore improves the controller and the estimator jointly. Mathematically, we solve the first stage with first-order methods (e.g.\ BFGS) driven by the stochastic optimal control gradient, which the envelope theorem supplies from the inner solver's dual variables (exactly under a unique inner minimizer, and as a subgradient otherwise). Across our studies this yields significant design cost reductions in which the sensing axis, absent from control-only co-design, is a major and sometimes dominant contributor.

\emph{Contributions:}
\begin{itemize}
    \item We formulate the control and estimation co-design as a single two-stage problem: a design first stage over a stochastic optimal control problem
    \item We show that the design gradients are obtained exactly from the dual variables, already returned by the solver. That is, no implicit differentiation of the Karush-Kuhn-Tucker (KKT) system is needed. Therefore, the gradient computation comes with no (or trivial) cost.
    \item We provide regularity conditions (a unique inner minimizer and strict complementarity) under which the envelope gradient is exact, and show that when they fail the same returned dual still yields a valid subgradient, so the method degrades, but does not fail.
    \item Casting the inner problems as SDPs buys modeling flexibility such as: variety of convex design regularizers, quadratic stability over a polytopic (convex-hull) uncertainty set, and structural/sparsity constraints on the covariances or the control and estimator gains.
    \item As an approximate extension to constrained nonlinear systems, a regime built around the linearization of eKF dynamics is used to construct a convex model predictive control (MPC) problem, with approximate gradients read from the
    dynamics costates (adjoint variables).
    \item Sensor allocation is absorbed directly into the formulation, either as a mixed-integer selection or as an $\ell_1$ relaxation promoting sparse sensing.
    \item We validate the framework on four realistic co-design studies across spacecraft, process, and power systems, with diverse objectives and modeling techniques.
\end{itemize}
\suggestnote{Consider adding a section pointer to each bullet (e.g. "(Sec.~3)", "(Sec.~4)"), the way many co-design/optimization papers do. A reader skimming the contributions list right now has no way to jump to where each claim is developed.}


%=======================================================================
\section{Related work}\label{sec:related}
%=======================================================================

\emph{Control co-design and integrated design.}
\cite{fathy2001coupling} established that plant and controller optimization are
coupled and cannot be separated without loss; \cite{allison2014multidisciplinary}
cast dynamic-system CCD as multidisciplinary design optimization (MDO); and
\cite{herber2019nested} systematized the nested and simultaneous solution
strategies. In all of these the estimator design and sensing layers are typically assumed fixed.

\emph{Structural co-design and information architecture.}
\cite{grigoriadis1998integrate} alternated a control linear matrix inequality (LMI) with a plant parameter design. Closest to the present work, \cite{li2008integrating} treated the information architecture (sensor/actuator gains) as decision variables jointly with control or estimation via LMIs, and \cite{saraf2017h2,das2017sparse} pushed this toward sparse $H_2$ sensing with guaranteed estimation-error bounds. These methods alternate LMIs over gains for a fixed plant and cost class, whereas ContEst places actuation, sensing, and the estimator itself in one stochastic objective and obtains the joint design gradient exactly rather than relying on alternation.

\emph{Sensor selection and LQG sensing co-design.}
A large literature selects sensors (or actuators) for estimation via
convex relaxation \cite{joshi2009sensor}, submodularity and greedy algorithms
\cite{tzoumas2016sensor,zhang2017sensor}, or sparsity-promoting selection
\cite{dhingra2014admm}. These covariance-based criteria are the control-theoretic
face of optimal experimental design (our estimation cost
$\operatorname{tr}(Q\Sigma^\varepsilon)$ is of A-optimal type, $\Sigma^\varepsilon$ is the state estimation error covariance), and the
identification community's least-costly experiment design
\cite{gevers2005identification} is a sibling of our noise level
co-design (Section~\ref{sec:weights}).
Closest to us, \cite{tzoumas2021lqg} co-design sensing with LQG control through
the separation structure (see also \cite{siami2021separation,zardini2021codesign});
but these choose \emph{which} sensors to deploy from a discrete menu, strictly in the
LQG regime.


\emph{Differentiable optimization and the envelope theorem.}
Recent work differentiates through convex control programs by implicitly differentiating the KKT system \cite{amos2017optnet,agrawal2019differentiable}. As we stress in
Section~\ref{sec:envelope}, for the optimal value (not the optimal argument) this machinery is unnecessary. Implicit differentiation of the KKT system solves an $m\times m$ linear system per design
direction, with $m=O(r_x^2)$ for the covariance SDPs ($r_x$ is the state-space dimension), i.e.\ $O(r_x^6)$ per
directional derivative, and it differentiates the minimizer, which ContEst
never needs. The envelope route reads the already-computed dual, therefore requires $O(r_x^2)$ for the whole gradient.

\emph{Parametric LQ and Riccati-based design gradients.}
The design gradient we recover is classical: the two-gramian formula dates to the
constant output feedback problem of \cite{levine1970determination} and the
parametric linear quadratic (LQ) machinery of \cite{makila1987computational,rautert1997computational}. On the estimation side \cite{belabbas2016geometric} differentiates the filter
Riccati, now a generic AD primitive \cite{kao2020autodiff} exploited by
adjoint-based co-design \cite{he2026efficient}. For unconstrained and plain linear quadratic regulator (LQR) inner problem, these
Riccati/adjoint routes are computationally cheaper than the SDPs we use, and ContEst can revert to the Riccati path when possible. Our
contribution is not the gradient in isolation but its unification, via the envelope theorem, across the control and estimation halves and their
$H_2$/$H_\infty$/eKF--MPC variants, and retention of the SDP path where it buys modeling power the Riccati route cannot express: quadratic stability over a convex uncertainty set, structural/sparsity
constraints, and others. Moreover, exploiting structure (diagonal or masked covariances) makes
the conic route cheaper in cost.

\emph{Positioning of the paper.}
Table~\ref{tab:positioning} gives a summary comparison of other methods per capability offered by ContEst (other methods may provide other benefits in other contexts). Nevertheless, ContEst need not operate in isolation; its only external output is a design gradient, it can be composed with the methods it is compared against rather than replacing them.

\begin{table*}[t]
\centering
\caption{Positioning of ContEst against the nearest lines of work, by capability
(\checkmark\ present, \ding{55}\ absent). ``Estimator
itself designed'' means tuning a dynamic estimator (e.g.\ bandwidth 
or noise level), not only sensor selection.}
\label{tab:positioning}
\explainnote{How to read this table: each column is a line of prior work, each row a capability; a checkmark means that line of work can do it, an $\times$ means it can't (or doesn't natively). The row that matters most for your pitch is "estimator itself designed", where ContEst is the only checkmark in the whole table -- that's the single strongest visual argument for novelty, worth calling out explicitly in the surrounding text (Sec.~2 already gestures at it but doesn't point at this specific row).}
\small
\resizebox{\textwidth}{!}{%
\begin{tabular}{lcccccc}
\toprule
capability & CCD/MDO & sensor & LQG sensing & info-architecture
& diff.\ optimization & \textbf{ContEst}\\
 & \cite{garcia2019control,herber2019nested} & selection
 \cite{joshi2009sensor,tzoumas2016sensor} & co-design \cite{tzoumas2021lqg}
 & LMIs \cite{li2008integrating}
 & layers \cite{amos2017optnet,agrawal2019differentiable} & \\
\midrule
plant-actuation co-designed & \checkmark & \ding{55} & \ding{55} & \checkmark & \ding{55} & \checkmark\\
plant-sensing (output map) co-designed & \ding{55} & \checkmark & \checkmark~(menu selection) & \checkmark & \ding{55} & \checkmark~(menu selection or $\ell_1$)\\
estimator itself designed & \ding{55} & \ding{55} & \ding{55} & \ding{55} & \ding{55} & \checkmark\\
closed-loop control in objective & \checkmark & \ding{55} & \checkmark & \checkmark & \checkmark & \checkmark\\
nonlinear regime & \checkmark & \ding{55} & \ding{55} & \ding{55} & \checkmark & \checkmark~(eKF-MPC surrogate)\\
gradient w.r.t.\ design & often FD & \ding{55} & \ding{55} & $\approx$ LMI alternation & needs KKT diff. & exact from duals\\
\bottomrule
\end{tabular}}
\end{table*}


Throughout, $\| x \|_Q^2 = x^\top Q x$, $\mathbb{E}$ denotes expectation,
$\operatorname{tr}(\cdot)$ the trace, and $M \succeq 0$ ($M \succ 0$) positive
(semi)definiteness. For a signal $s$, $s_{i:j} = (s_i, s_{i+1}, \dots, s_j)$ denotes
the sequence from time $i$ to $j$. The information available at time $k$. A superscript ${}^\star$ marks an optimal value.
\errornote{"The information available at time $k$." is not a complete sentence -- it looks like a leftover fragment. Likely missing subject/verb, e.g. "$\mathcal Z_k$ denotes the information available at time $k$." Please fix before submission.}

%=======================================================================
\section{Problem formulation}\label{sec:formulation}
%=======================================================================

Consider the parametrized stochastic system
\begin{subequations}\label{eq:ss}
\begin{align}
x_{k+1} &= f_\theta(x_k, u_k) + w_k, \label{eq:ss-state}\\
y_k     &= h_\theta(x_k) + v_k, \label{eq:ss-output}
\end{align}
\end{subequations}
with state $x_k \in \mathbb{R}^{r_x}$, input $u_k \in \mathbb{R}^{r_u}$, and
output $y_k \in \mathbb{R}^{r_y}$. The maps $f_\theta, h_\theta$ are parametrized by $\theta \in \mathbb{R}^{r_\theta}$, representing plant parameters, actuator gains, and sensor tuning/placement. The state dynamics $f_\theta$ is locally bounded in its arguments for all $\theta$. The disturbances $w_k, v_k$ are i.i.d.,
mutually independent and independent of $x_0 \sim p_0$ (the prior, before any
measurement including $y_0$), with covariances $W,V$, respectively, which are possibly parametrized by $\theta$ as well.

The goal is to jointly optimize design, control, and estimation altogether. The
design cost $J_{\mathrm{des}}(\theta)$ and design set $\theta \in \Theta$
encode the price, manufacturability, placement difficulty or some mechanical, thermal, structural, and other considerations. The
control and estimation costs are embedded in the stochastic optimal control problem
\begin{equation} \label{eq:Jstoch_N}
\begin{aligned}
J_{\mathrm{sc}}^\star(\mathcal Z_0;\theta) &= \min_{u_{0:N-1}} J_{\mathrm{sc}}(\mathcal Z_0,u_{0:N-1};\theta),\\
&=\min_{u_{0:N-1}} \frac{1}{N}\,
\mathbb{E}\big\{ \sum_{k=0}^{N-1} \ell(x_k,u_k) + \ell^N(x_N) \big\},
\end{aligned}
\end{equation}
for finite-horizon problems, or for infinite-horizon as in 
\begin{equation}\label{eq:Jstoch_infty}
J_{\mathrm{sc}}^\star(\mathcal Z_0;\theta) = \min_{\kappa} \lim_{k \to \infty} \mathbb{E} \, \ell(x_k,\kappa(\mathcal Z_k)),
\end{equation}
with the $\ell, \ell^N$ are stage and terminal costs are typically
\begin{equation*}
   \ell(x_k,u_k) = \|x_k \|_Q ^2 + \| u_k\|_R^2, \quad \ell_N(x_N) = \| x_N \|_Q^2,
\end{equation*}
\errornote{Terminal cost notation is inconsistent: it's $\ell^N(x_N)$ (superscript) in Eq.~(2) above but $\ell_N(x_N)$ (subscript) here. Pick one and use it everywhere -- I'd suggest $\ell^N$ to avoid clashing with the time-index subscripts used everywhere else.}
and the information up to time $k$, $\mathcal Z_k = (p_0, y_{0:k}, u_{0:k-1}$ and $\mathcal Z_0 = (p_0, y_0)$.
\errornote{Missing closing parenthesis: should read $\mathcal Z_k = (p_0, y_{0:k}, u_{0:k-1})$ -- as written the paren opened before $y_{0:k}$ is never closed.}

The above minimizations are possibly subject to state and input constraints and the expectation is over $(x_0, w_{0:N-1}, v_{0:N})$.

\begin{prob}[ContEst]\footnote{Problem~\ref{prob:bilevel} is a two-stage problem, a special form of a bilevel problem; the outer optimization contains the value of the lower-level problem, not its argmin. Solving a bilevel problem directly is hard since it requires implicitly differentiating the inner optimality conditions. Instead, in the two-stage case, the envelope theorem supplies the gradient
directly from the duals the inner solver already returns.}\label{prob:bilevel}
\begin{align*}
&\text{Design (first stage):}\quad \min_{\theta \in \Theta}\
J(\theta) = J_{\mathrm{des}}(\theta) + J_{\mathrm{sc}}^\star(\mathcal Z_0;\theta),\\
&\text{Control (second stage):}\quad J_{\mathrm{sc}}^\star(\mathcal Z_0;\theta)\ \text{as in
\eqref{eq:Jstoch_N} or \eqref{eq:Jstoch_infty}}.
\end{align*}
\end{prob}

\begin{rem}[Estimation co-design]
    Equations \eqref{eq:Jstoch_N} and \eqref{eq:Jstoch_infty}, although appear to be optimizing over control only, they implicitly require optimal estimation as well. The control policies are feedback in the information state (state filtered density), which require careful filter design for their reduction, approximation and accuratel tracking \cite{kumar2015stochastic}.
\errornote{Typo: "accuratel" $\to$ "accurate".}
\end{rem}

\begin{rem}[Operating scenarios]\label{rem:scenarios}
Problem~\ref{prob:bilevel} is stated for a single prior $p_0$. Multiple operating scenarios: initial densities, tracking trajectories/set-point, disturbance statistics, can be accommodated by replacing the second-stage value with the, say, a weighted average of the cost over the different scenarios.
\end{rem}

Note that the second stage is a stochastic (dual) control problem whose sufficient
statistic is the information state: the filtered density in this paper. The exact reduction of this problem into two LQG SDPs in the next section, and approximate reduction to eKF-MPC in a later section are detailed in Appendix~\ref{app:infostate}.

%=======================================================================
\section{The LQG case: LQR and Kalman SDPs}\label{sec:lqg}
%=======================================================================

For a linear (or linearized) plant
$$x_{k+1}=A_\theta x_k+B_\theta u_k+w_k, \quad y_k=C_\theta x_k+v_k,$$ with zero-mean
the inner problem is an LQG problem whose cost splits, by
the separation principle, into control and estimation parts,
\add{with $Q \succeq 0 \in \mathbb{R}^{r_x \times r_x}$ and
$R \succ 0 \in \mathbb{R}^{r_u \times r_u}$ the state and input cost
weights already used in the stage cost $\ell$ of
Section~\ref{sec:formulation} (Proposition~\ref{prop:lqgreg} below
assumes the stronger $Q \succ 0$), and control input
$u_k = Kx_{k \mid k}$ the certainty-equivalent state-feedback law
(justified by the separation principle, see Appendix~\ref{app:infostate};
$K$ is a free gain here, not yet the optimal one -- it becomes the
optimal LQR gain once \eqref{eq:Jcont} below is solved)}
\begin{equation} \label{eq:linear stoch cost}
\begin{aligned}
&J_{\mathrm{sc}}=J_{\mathrm{est}}(\theta)+J_{\mathrm{cont}}(\theta),\\
&J_{\mathrm{est}}=\operatorname{tr}(Q\Sigma^\varepsilon), \quad
J_{\mathrm{cont}}=\operatorname{tr}(Q\Sigma)+\operatorname{tr}(RK\Sigma K^\top)
\end{aligned}
\end{equation}
(Appendix~\ref{app:infostate} contains the full derivation), where $\Sigma=\lim_{k \to \infty}\mathbf{Cov}( x_{k \mid k})$ is the
steady-state estimated state covariance ($x_{k \mid k}$ is the state estimate given $\mathcal{Z}_k$), and $\Sigma^\varepsilon=\lim_{k \to \infty}\mathbf{Cov}(\varepsilon_k)$ the steady-state
estimation-error covariance. \add{With $G$ the observer (Kalman) gain,} the observer is given by
\begin{align*}
    x_{k+1 \mid k+1} = A_\theta x_{k \mid k} + B_\theta u_k + G(y_k - C_\theta x_{k \mid k}).
\end{align*}
\errornote{Both $K$ (introduced only implicitly, just above in Eq.~\eqref{eq:linear stoch cost}) and $G$ (here) get used before ever being properly defined in the main text. $K$'s actual definition -- $u_k=\kappa(\mathcal Z_k)=Kx_{k\mid k}$, certainty-equivalent state feedback, justified by the separation principle -- only appears in Appendix~\ref{app:infostate}, which comes \emph{after} this section in reading order. $G$ never gets an explicit defining sentence anywhere in the paper at all; its role as the observer/innovation gain is only inferable from the structure of this one equation. I've added short inline reminders at both spots (in red), but the cleaner fix would be to move the "$u_k=Kx_{k\mid k}$, certainty equivalent control" sentence from Appendix~\ref{app:infostate} up to right before Eq.~\eqref{eq:linear stoch cost}, so both gains are named before they're used rather than after. Same recurring pattern as the $Q$ gap at Eq.~(4) and the $W,V$ gap at Eq.~(A.4)/(A.5) flagged earlier: symbols defined far from, or after, their first use.}
For a
fixed $\theta$, the two parts in \eqref{eq:linear stoch cost} depend on disjoint variables and are minimized
independently \cite{tzoumas2021lqg}, but at the design level $\theta$ couples
$A_\theta,B_\theta,C_\theta$ and the noise covariances, and hence, both the control
gain $K$ and the observer gain $G$. Therefore, co-design is not separable across the two axes, which is precisely why the joint
two-stage formulation of Problem~\ref{prob:bilevel} is needed rather than a
sequential pipeline.

\emph{Control SDP.} Following the linearizing change of
variables $L=K\Sigma$ (with the epigraph slack $Z_0$ upper-bounding the input
second moment $K\Sigma K^\top$ through the first LMI) convexifies the regulator
into the semidefinite program
\begin{equation}\label{eq:Jcont}
\begin{aligned}
J_{\mathrm{cont}}^\star(\theta) = &\min_{\Sigma,L,Z_0}\
\operatorname{tr}(Q\Sigma) + \operatorname{tr}(R Z_0)\\
\text{s.t. } & \Sigma \succ 0, \quad
\begin{bmatrix} Z_0 & L \\ L^\top & \Sigma \end{bmatrix}
\succeq 0,\\
& \begin{bmatrix}
\Sigma - W & A_\theta\Sigma + B_\theta L \\
(A_\theta\Sigma + B_\theta L)^\top & \Sigma
\end{bmatrix} \succeq 0,
\end{aligned}
\end{equation}
with optimal gain $K^\star = L^\star (\Sigma^\star)^{-1}$.

\emph{Estimation SDP.} Dually, and moving from the controllability gramian $\Sigma^\varepsilon$ to the observability one $P^\varepsilon$ that solves a
Lyapunov inequality, which after the change of variables $F = P^\varepsilon G$
(with slack $Z_1$) and the LMI relaxation, gives
\begin{equation}\label{eq:Jest}
\begin{aligned}
J_{\mathrm{est}}^\star(\theta) = &\min_{P^\varepsilon,F,Z_1}\
\operatorname{tr}(W P^\varepsilon) + \operatorname{tr}(V Z_1)\\
\text{s.t. } & P^\varepsilon \succ 0, \quad
\begin{bmatrix} Z_1 & F^\top \\ F & P^\varepsilon \end{bmatrix}
 \succeq 0,\\
& \begin{bmatrix}
P^\varepsilon - Q & A_\theta^\top P^\varepsilon - C_\theta^\top F^\top \\
P^\varepsilon A_\theta - F C_\theta & P^\varepsilon
\end{bmatrix} \succeq 0,
\end{aligned}
\end{equation}
dual in form to \eqref{eq:Jcont} and exposing the sensing operator $C_\theta$,
i.e.\ the estimation half of the co-design \cite{li2008integrating}.

\emph{Design gradients from the LMI duals.} Both \eqref{eq:Jcont} and
\eqref{eq:Jest} are instances of the generic conic program, so
the envelope theorem (will be detailed in Section~\ref{sec:envelope}) delivers their design gradients
directly from the returned LMI duals, with no differentiation of the solver.
Applied to the Lyapunov constraint of \eqref{eq:Jcont},
whose dual we partition as
$S = \left[\begin{smallmatrix} S_{11} & S_{12} \\ S_{12}^\top & S_{22}
\end{smallmatrix}\right] \succeq 0$, gives (concretely, the envelope theorem
evaluates $\partial L/\partial\theta$ of the Lagrangian at the optimal primal
and dual values already returned by the solver, so no constraint need be
re-differentiated)
\begin{equation}\label{eq:gradAB}
\frac{\partial J_{\mathrm{cont}}^\star}{\partial A_\theta} = -2\, S_{12}^\star\Sigma^\star,
\quad
\frac{\partial J_{\mathrm{cont}}^\star}{\partial B_\theta} = -2\, S_{12}^\star(L^\star)^\top,
\end{equation}
Partitioning the dual of the estimation LMI in \eqref{eq:Jest} the same way,
$S' = \left[\begin{smallmatrix} S'_{11} & S'_{12} \\ S'^\top_{12} & S'_{22}
\end{smallmatrix}\right] \succeq 0$, its off-diagonal block $P^\varepsilon A_\theta
- F C_\theta$ yields the parallel pair
\begin{equation}\label{eq:gradAC}
\frac{\partial J_{\mathrm{est}}^\star}{\partial A_\theta} = 2\, P^{\varepsilon\star} S^{\prime\star}_{12},
\quad
\frac{\partial J_{\mathrm{est}}^\star}{\partial C_\theta} = -2\, F^{\star} S^{\prime \star}_{12}.
\end{equation}
\explainnote{Worth a one-line remark here on why the operand order flips between~\eqref{eq:gradAB} ($S_{12}^\star\Sigma^\star$, dual-then-primal) and~\eqref{eq:gradAC} ($P^{\varepsilon\star}S_{12}^{\prime\star}$, primal-then-dual): it follows from the control/estimation problems being transposed duals of one another, but a reader seeing these side by side for the first time will likely wonder if it's a typo. Also double check the sign: \eqref{eq:gradAB} is $-2S_{12}^\star\Sigma^\star$ (minus) while \eqref{eq:gradAC} is $+2P^{\varepsilon\star}S_{12}^{\prime\star}$ (plus) -- please confirm this sign flip is intentional (it plausibly is, from the transpose/duality map), not a slip.}
The design gradient follows by the chain rule
(Section~\ref{sec:envelope}),
\begin{align}\label{eq:chain}
\frac{\partial J_{\mathrm{cont}}^\star}{\partial \theta_i} &=
\operatorname{tr}\!\Big(\frac{\partial J_{\mathrm{cont}}^\star}{\partial A_\theta}^{\!\top}
\frac{\partial A_\theta}{\partial \theta_i}\Big)
+ \operatorname{tr}\!\Big(\frac{\partial J_{\mathrm{cont}}^\star}{\partial B_\theta}^{\!\top}
\frac{\partial B_\theta}{\partial \theta_i}\Big) ,\nonumber\\
&= \left \langle\frac{\partial J_{\mathrm{cont}}^\star}{\partial A_\theta},
\frac{\partial A_\theta}{\partial \theta_i}\right \rangle
+ \left \langle\frac{\partial J_{\mathrm{cont}}^\star}{\partial B_\theta},
\frac{\partial B_\theta}{\partial \theta_i}\right \rangle,
\end{align}
and the gradient vector is constructed from its components
\begin{align*}
    \nabla_\theta J_{\mathrm{sc}}^\star = \Big [\left(\frac{\partial J_{\mathrm{cont}}^\star}{\partial \theta_1} + \frac{\partial J_{\mathrm{est}}^\star}{\partial \theta_1}\right )\,\, \left (\frac{\partial J_{\mathrm{cont}}^\star}{\partial \theta_2} + \frac{\partial J_{\mathrm{est}}^\star}{\partial \theta_2} \right )\, \cdots  \Big]^\top.
\end{align*}


\subsection{Sensor precision: noise level co-design}\label{sec:weights}
The weights $Q, R$ of the quadratic cost and the noise covariances $W, V$ of the
process and measurement can also be parametrized by $\theta$. This matters in practice because they can represent the knobs a sensor or actuator budget buys.

Concretely, both LQG subproblems are affine in these data, so their sensitivities
are read off directly, with no LMI-dual partition required for the objective
terms. From the control SDP~\eqref{eq:Jcont},
\begin{equation}\label{eq:gradQRcont}
\frac{\partial J_{\mathrm{cont}}^\star}{\partial Q_\theta} = \Sigma^\star,\quad
\frac{\partial J_{\mathrm{cont}}^\star}{\partial R_\theta} = Z_0^\star,\quad
\frac{\partial J_{\mathrm{cont}}^\star}{\partial W_\theta} = -\,S_{11}^\star,
\end{equation}
the first two from the objective $\operatorname{tr}(Q_\theta\Sigma)+\operatorname{tr}(R_\theta Z_0)$ and the
last from the $(1,1)$ block of $S^\star$. From the estimation SDP~\eqref{eq:Jest},
\begin{equation}\label{eq:gradQRest}
\frac{\partial J_{\mathrm{est}}^\star}{\partial W_\theta} = P^{\varepsilon\star},\quad
\frac{\partial J_{\mathrm{est}}^\star}{\partial V_\theta} = Z_1^\star,\quad
\frac{\partial J_{\mathrm{est}}^\star}{\partial Q_\theta} = -\,S_{11}^{\prime\star},
\end{equation}
with $S_{11}^{\prime\star}$ the $(1,1)$ dual block of the estimation LMI. The total gradient is then the sum of all the contributions, added to \eqref{eq:chain} via chain rule.

\begin{rem}[Modeling flexibility]
Posing both halves as SDPs is not only convenient for the dual-based gradients:
it also opens the formulation to constraints a plain Riccati recursion cannot
express: LMI regularizers (e.g.\ trace or nuclear-norm penalties promoting sparse
actuation/sensing), budget and $\ell_1$ sensor-selection constraints on $\theta$,
or quadratic stability robustness constraints over a polytopic data set. Proposition~\ref{prop:regexact} in the next section
shows every such convex ingredient is inherited by the exact-gradient machinery of
Section~\ref{sec:envelope}.
\end{rem}

%=======================================================================
\section{Exact design gradients via the envelope theorem}\label{sec:envelope}
%=======================================================================

The two-stage problem is solved by descent on the first-stage variable $\theta$,
so the crux is $\nabla_\theta J_{\mathrm{sc}}^\star$: the sensitivity of a
\emph{value function} to parameters entering its objective and constraints. Our
central observation is that this gradient is essentially free---and it is exactly
this that reduces the bilevel program to the two-stage one. Write the convex
second-stage problem in conic form,
\begin{equation}\label{eq:generic}
V(\theta) = \min_{z}\ \phi(z,\theta)\quad \text{s.t.}\quad g(z,\theta)=0,\ \
G(z,\theta) \in \mathcal{K},
\end{equation}
with closed convex cone $\mathcal{K}$ (the positive semidefinite cone for SDPs),
Lagrangian $L(z,\mu,S,\theta) = \phi(z,\theta) + \langle \mu, g(z,\theta)\rangle
+ \langle S, G(z,\theta)\rangle$, and optimal primal--dual triple
$(z^\star,\mu^\star,S^\star)$. We collect the hypotheses under which the value
$V$ is differentiable and its gradient is read from that triple.

\begin{assumption}\label{as:env}
On a neighborhood $\Theta_0$ of the design point $\theta$:
\emph{(A1)}~for each $\theta\in\Theta_0$, \eqref{eq:generic} is convex in $z$,
its data $(\phi,g,G)$ are jointly $C^1$ in $(z,\theta)$, and Slater's condition
holds, so strong duality holds with an attained primal--dual solution;
\emph{(A2)}~the primal minimizer $z^\star(\theta)$ is unique;
\emph{(A3)}~the primal--dual solution is \emph{nondegenerate} and strictly
complementary, so the dual $(\mu^\star,S^\star)$ is unique
\cite{alizadeh1997complementarity,shapiro1997differentiability};
\emph{(A4)}~$\theta$ enters \eqref{eq:generic} only through finitely many data
objects $D(\theta)$ (e.g.\ $A_\theta,B_\theta,C_\theta,W,V$) that are $C^1$
in $\theta$.
\end{assumption}
\explainnote{Plain-language gloss on Assumption~\ref{as:env}, for readers who don't work in conic programming daily: (A1) says the inner problem is a "nice" convex program with no duality gap; (A2) says there's exactly one best solution, not a tie; (A3) says the same holds for the dual/multiplier side (this is the one that fails at degenerate optima, e.g.\ Remark~\ref{rem:exactfail}); (A4) says $\theta$ only shows up through a handful of matrices, so the chain rule in Section~\ref{sec:envelope}'s point (3) is cheap. Might be worth adding a sentence like this right after the assumption for accessibility.}

Assumptions (A1)--(A3) are exactly the regularity under which the optimal value
of a conic (in particular semidefinite) program is differentiable and its
sensitivity is governed by the unique dual
\cite{bonnans2000perturbation,shapiro1997differentiability}; (A4) is a modeling
convenience that isolates the $\theta$-dependence in cheap model Jacobians.

\begin{theorem}[Exact design gradient from the inner dual]\label{prop:envelope}
Under Assumption~\ref{as:env}, $V$ is differentiable at $\theta$ and
\begin{equation}\label{eq:envelope}
\nabla_\theta V(\theta) = \nabla_\theta L(z^\star,\mu^\star,S^\star,\theta)
= \partial_\theta \phi + \langle \mu^\star, \partial_\theta g\rangle
+ \langle S^\star, \partial_\theta G\rangle,
\end{equation}
all partials evaluated at $(z^\star,\theta)$ with $z^\star,\mu^\star,S^\star$
held fixed.
\end{theorem}
\begin{proof}
By (A1)--(A3) the solution map $\theta\mapsto(z^\star,\mu^\star,S^\star)$ is
single-valued and differentiable on $\Theta_0$
\cite{bonnans2000perturbation,shapiro1997differentiability}, and strong duality
with complementary slackness gives
$V(\theta) = L(z^\star(\theta),\mu^\star(\theta),S^\star(\theta),\theta)$.
Differentiating and using inner stationarity $\nabla_z L = 0$ and the
stationarity of $L$ in $(\mu,S)$ at the solution,
$\nabla_\theta V = \nabla_z L\,\partial_\theta z^\star
+ \nabla_{(\mu,S)}L\,\partial_\theta(\mu^\star,S^\star) + \partial_\theta L
= \partial_\theta L$ \cite{milgrom2002envelope,bonnans2000perturbation}.
\end{proof}
\explainnote{Proof in one line, for the reader: differentiate $V(\theta)=L(z^\star(\theta),\ldots,\theta)$ by the chain rule; the terms tracking how $z^\star,\mu^\star,S^\star$ move with $\theta$ all vanish because the inner problem's own first-order (stationarity) conditions say $\partial L/\partial z=0$, $\partial L/\partial \mu = 0$, $\partial L/\partial S=0$ at the optimum -- so only the "direct" $\partial_\theta L$ term survives. This is the whole reason the gradient is free: nobody needs to compute $\partial_\theta z^\star$.}

\begin{proposition}[Exact gradient preserved under added convex structure]\label{prop:regexact}
Augment \eqref{eq:generic} with any finite family of additional convex,
conic-representable ingredients: regularizers $r_i(z,\theta)$ added to the objective
$\phi$, and/or constraints $G_i(z,\theta)\in\mathcal{K}_i$ with closed convex cones
$\mathcal{K}_i$, all data $C^1$ in $(z,\theta)$. If the augmented program still
satisfies Assumption~\ref{as:env}, then $V$ remains differentiable at $\theta$ and
\begin{equation}\label{eq:envelope_aug}
\nabla_\theta V(\theta)=\partial_\theta\phi+\sum_i\partial_\theta r_i
+\langle\mu^\star,\partial_\theta g\rangle+\langle S^\star,\partial_\theta G\rangle
+\sum_i\langle\lambda_i^\star,\partial_\theta G_i\rangle,
\end{equation}
i.e.\ each added ingredient contributes exactly one envelope term of the same dual
form, with no re-differentiation. A sufficient recipe for the two binding hypotheses
is: \emph{(A2)} holds whenever some added regularizer is \emph{strictly convex}
(e.g.\ a Tikhonov $\tfrac{\rho}{2}\|z\|^2$ or a $-\log\det$ barrier), forcing a
unique minimizer; \emph{(A3)} holds whenever the active constraints have linearly
independent supports and are strictly complementary (automatic for inactive ones)
\cite{alizadeh1997complementarity,shapiro1997differentiability}.
\end{proposition}
\begin{proof}
Absorb the extra cones into the product cone
$\mathcal{K}\times\prod_i\mathcal{K}_i$ and the extra data into $D(\theta)$: the
augmented program is again an instance of \eqref{eq:generic}, so
Theorem~\ref{prop:envelope} applies verbatim. Its dual term
$\langle S^\star,\partial_\theta G\rangle$ splits into the base block plus one block
$\langle\lambda_i^\star,\partial_\theta G_i\rangle$ per added constraint, and the
strictly-convex regularizer partials appear in $\partial_\theta\phi$, giving
\eqref{eq:envelope_aug}. The (A2)/(A3) recipe is standard: strict convexity yields a
unique $z^\star$, and nondegenerate strictly-complementary active supports yield a
unique dual \cite{alizadeh1997complementarity,shapiro1997differentiability}.
\end{proof}

\begin{rem}[Where exactness is lost]\label{rem:exactfail}
The recipe of Proposition~\ref{prop:regexact} fails exactly when (A3) does: at a
\emph{degenerate active face} (rank-deficient active supports), or when a constraint
\emph{binds at its achievable floor} (loss of strict complementarity---for instance
the minimal-$\gamma$ optimum of Section~\ref{sec:hinf}). The optimal dual is then
set-valued and the exact identity \eqref{eq:envelope_aug} need not hold; the next
proposition shows the framework degrades gracefully rather than breaking. In our
studies strict complementarity held at every solution (consistent with the gradient
check of Section~\ref{sec:numerical}), and by Proposition~\ref{prop:lqgreg} the LQG
regime can degenerate only where its hypotheses fail (singular $W$, or lost
detectability at isolated $\theta$), pinpointing in advance the points that need the
fallback.
\end{rem}

\begin{proposition}[Subgradient fallback under (A1)]\label{prop:subgrad}
Suppose (A1) and (A4) hold on a neighborhood of $\theta$ but (A2) or (A3) fails.
Then $V$ is locally Lipschitz and directionally differentiable, with
\begin{equation}\label{eq:dirderiv}
V'(\theta;d)=\min_{(\mu,S)\in\mathcal{D}^\star(\theta)}
\langle\nabla_\theta L(z^\star,\mu,S,\theta),d\rangle,
\end{equation}
where $\mathcal{D}^\star(\theta)$ is the compact convex set of optimal duals; and the
particular dual $(\mu^\star,S^\star)$ the inner solver returns is a genuine
subgradient, $\partial_\theta L(z^\star,\mu^\star,S^\star,\theta)\in\partial V(\theta)$
(the Clarke subdifferential). Hence $\partial_\theta L$ supplies a descent direction,
first-order stationarity reads $0\in\partial V(\theta)$, and the subgradient/proximal
outer method of Remark~\ref{rem:algorithm} converges---at a linear rather than
superlinear rate, exactly as nonsmooth $H_\infty$ synthesis relies on
\cite{bonnans2000perturbation}. Where (A1) itself fails (a feasibility boundary with
$V\to+\infty$, or loss of Slater/dual boundedness) no finite subgradient need exist.
\end{proposition}
\begin{proof}
Under (A1) the value $V$ is locally Lipschitz---Appendix~\ref{app:subgrad} gives the
quantitative bound, from dual feasibility being $\theta$-independent and $L$ being
$C^1$ with bounded optimal primal/dual sets. By Rademacher's theorem $V$ is
differentiable a.e., and at each such point $\nabla V=\partial_\theta L$ by
Theorem~\ref{prop:envelope}; passing to the limit places
$\partial_\theta L(z^\star,\mu^\star,S^\star,\theta)$ in the nonempty compact convex
Clarke set $\partial V(\theta)$, which collapses to the singleton
$\{\nabla V(\theta)\}$ precisely when (A2)--(A3) also hold. The directional-derivative
formula \eqref{eq:dirderiv} is the standard conic-perturbation result
\cite{bonnans2000perturbation}.
\end{proof}

Three points give ContEst its practical edge. \emph{(1)~No implicit
differentiation:} \eqref{eq:envelope} has no term $\partial_\theta z^\star$; we
never compute how the minimizer moves with $\theta$, which is precisely what
implicit-function-theorem methods spend effort on
\cite{amos2017optnet,agrawal2019differentiable}. \emph{(2)~No KKT
differentiation:} the multipliers $(\mu^\star,S^\star)$ are not recomputed by
differentiating optimality conditions; they are the dual solution the inner
solver already returns, and we read them off. \emph{(3)~Only model Jacobians
remain:} in our setting $\theta$ enters \eqref{eq:generic} through a small set of
data objects $D$ (e.g.\ the matrices $A_\theta, B_\theta, \dots$), so
$\nabla_\theta V = \sum_D \langle \partial V / \partial D,\ \partial D /
\partial\theta\rangle$, where $\partial V/\partial D$ comes from
\eqref{eq:envelope} and $\partial D/\partial\theta$ are explicit (or
automatically differentiated) Jacobians of the \emph{model}, not the optimizer.

We now verify, for the LQG SDPs of Section~\ref{sec:lqg}, the regularity of
Assumption~\ref{as:env} from the model data.

\begin{proposition}[LQG regularity]\label{prop:lqgreg}
Fix $\theta$ and suppose $(A_\theta,B_\theta)$ is stabilizable, $(A_\theta,
C_\theta)$ is detectable, and $Q,R,W,V\succ0$, with all data $C^1$ in $\theta$.
Then for the control SDP~\eqref{eq:Jcont} and the estimation
SDP~\eqref{eq:Jest}: \emph{(a)}~Slater's condition holds; \emph{(b)}~the primal
minimizer is unique---$\Sigma^\star$ the stabilizing closed-loop covariance and
$K^\star=L^\star(\Sigma^\star)^{-1}$ the Riccati gain, and dually
$P^{\varepsilon\star}=\Sigma^\varepsilon$ the Kalman error covariance;
\emph{(c)}~the Lyapunov-LMI dual is unique and strictly complementary---in fact
the $(1,1)$ dual block of the control LMI \textbf{is} the control Riccati
solution $P$ (so $S_{11}^\star=P$) and the $(1,1)$ dual block of the estimation
LMI \textbf{is} the steady-state Kalman error covariance $\Sigma^\varepsilon$ (so
$S_{11}^{\prime\star}=\Sigma^\varepsilon$). Hence Assumption~\ref{as:env} holds
on a neighborhood $\Theta_0$ of $\theta$, and Corollary~\ref{cor:lqg} applies.
\end{proposition}

\emph{Proof sketch.} A strictly feasible point is built from any stabilizing gain
inflated by $\varepsilon I$ (Slater); with $W\succ0$ a Lyapunov-eigenvector
argument forces every feasible $\Sigma$ to encode a stabilizing gain, and the
completion-of-squares identity $J(K)-\operatorname{tr}(PW)=\operatorname{tr}\!\big(
(K-K^\star)^\top(R+B^\top PB)(K-K^\star)\Sigma_K\big)$ gives primal uniqueness;
complementary slackness reduces the dual stationarity to the closed-loop Stein
equation $\Lambda=Q+K^{\star\top}RK^\star+A_{\mathrm{cl}}^\top\Lambda
A_{\mathrm{cl}}$, whose unique solution is $\Lambda=P$, and strict
complementarity follows from rank counting with $P\succeq Q\succ0$, $R\succ0$.
The full argument and the dual identities are in Appendix~\ref{app:lqgreg}.
\hfill$\square$

Controllability and observability imply the stabilizability/detectability
hypotheses, so the practical reading of Proposition~\ref{prop:lqgreg} is:
\emph{controllable, observable plant with positive-definite weights $\Rightarrow$
everything in Sections~\ref{sec:lqg} and~\ref{sec:envelope} is exact and
differentiable}. It also explains structurally why the Riccati gradient formulas
of Section~\ref{sec:sim:complexity} coincide with the LMI-dual
formulas~\eqref{eq:gradQRcont}--\eqref{eq:gradQRest}: they are literally the same
object, $S_{11}^\star=P$.

\begin{corollary}[LQG design gradient]\label{cor:lqg}
Under the hypotheses of Proposition~\ref{prop:lqgreg} (stabilizability,
detectability, $Q,R,W,V\succ0$), the SDPs~\eqref{eq:Jcont}--\eqref{eq:Jest}
satisfy Assumption~\ref{as:env}, so the Lyapunov-LMI dual $S$ (resp.\ its
estimation counterpart) is unique. Then $J_{\mathrm{cont}}^\star(\theta)$ and
$J_{\mathrm{est}}^\star(\theta)$ are differentiable at $\theta$ and their design
gradients are given exactly by \eqref{eq:gradAB}--\eqref{eq:chain} (resp.\ the
parallel expressions in $C_\theta, W, V$), read from the returned LMI duals.
\end{corollary}


%=======================================================================
\section{The $H_\infty$ analog}\label{sec:hinf}
%=======================================================================

The LQG development of Section~\ref{sec:lqg} considers average performance
`against' stochastic white disturbances. When the disturbances are better modeled as
unknown-but-bounded $\ell_2$ signals and the designer cares about the
worst case, an $H_\infty$ approach can be used. The inner value is now the
squared worst-case closed-loop $\ell_2$-gain $\gamma^2(\theta)$ rather than a
stationary variance.

Write the linear plant with an explicit disturbance channel and a performance
output,
\begin{equation}\label{eq:hinf_plant}
x_{k+1} = A_\theta x_k + B_\theta u_k + E_\theta w_k,\quad
z_k = C_{z}\,x_k + D_{z}\,u_k,
\end{equation}
where $w_k\in\ell_2$ is the exogenous disturbance, $E_\theta$ its input map
(e.g.\ $E_\theta = W_\theta^{1/2}$, so the process noise level becomes a disturbance weight), and $z_k$ stacks the penalized
state and input, $C_z = [\,Q^{1/2}\;\,0_{r_x \times r_u}\,]^\top$, $D_z = [\,0_{r_u \times r_x}\;\,R^{1/2}\,]^\top$, so the
$H_\infty$ objective weighs exactly the same signals as in $J_{\mathrm{sc}}^\star$.

\emph{Control SDP.} By the discrete-time bounded-real lemma (BRL) with the standard
linearizing change of variables $Y = P^{-1}\succ 0$, $L = KY$
\cite{boyd1994linear,grigoriadis1998integrate}, the state-feedback $H_\infty$
synthesis problem is the SDP (with $\mu := \gamma^2$)
\begin{equation}\label{eq:Hinf_cont}
\begin{aligned}
&\gamma^2_{\mathrm{cont}}(\theta) = \min_{Y,L,\mu}\ \mu \\
&\text{s.t. }  Y \succ 0,\ \
\begin{bmatrix}
-Y & \star & \star & \star \\
0 & -\mu I & \star & \star \\
A_\theta Y + B_\theta L & E_\theta & -Y & \star \\
C_z Y + D_z L & 0 & 0 & -I
\end{bmatrix} \preceq 0,
\end{aligned}
\end{equation}
(the $\star$ blocks are inferred by symmetry) with optimal gain
$K^\star = L^\star (Y^\star)^{-1}$.

\emph{Design gradients from the duals.} \eqref{eq:Hinf_cont} is a generic SDP
instance: the constraint is a single conic inclusion. Partition its dual
conformally with the $4\times4$ LMI,
$S_{\infty} = [\,S_{\infty,ij}\,]$, $S_{\infty,ij}=S_{\infty,ji}^\top$. The plant blocks
$A_\theta Y + B_\theta L$ sit in position $(3,1)$,
\begin{equation}\label{eq:gradHinf}
\frac{\partial \gamma^{\star\,2}_{\mathrm{cont}}}{\partial A_\theta} = 2\,S_{\infty,31}^\star\,Y^\star,
\quad
\frac{\partial \gamma^{\star\,2}_{\mathrm{cont}}}{\partial B_\theta} = 2\,S_{\infty,31}^\star\,(L^\star)^\top,
\end{equation}
with analogous expressions for $E_\theta$ (from the $(3,2)$ block $S_{\infty,32}^\star$),
and the design gradient assembled by the same chain rule~\eqref{eq:chain}.

\begin{corollary}[$H_\infty$ minimal-$\gamma$ design: a subgradient]\label{cor:hinf}
The expression \eqref{eq:gradHinf} for the minimal-$\gamma$ value
$\gamma^{\star\,2}_{\mathrm{cont}}(\theta)$ is, in general, only a subgradient: it
supplies a descent direction and the stationarity certificate
$0\in\partial\gamma^{\star\,2}_{\mathrm{cont}}$, but need not equal the gradient. Fixing $\gamma$ instead restores exactness, as we show next.
\end{corollary}
\begin{proof}
At the optimum \eqref{eq:Hinf_cont} drives $\mu$ down to the optimal
attenuation $\gamma^{\star\,2}_{\mathrm{cont}}$, where the BRL LMI loses rank: the $(2,2)$
block $-\mu I$ tightens against the Schur complement of the plant blocks, so
the active face is degenerate and strict complementarity (A3) generically
fails. By Proposition~\ref{prop:subgrad} the value is then locally Lipschitz
and the BRL dual the solver returns is a genuine subgradient of
$\gamma^{\star \, 2}_{\mathrm{cont}}$.
\end{proof}
\explainnote{In plain terms: at the best-possible attenuation $\gamma^\star$, the LMI is "just barely" tight and loses rank there, which is exactly the degenerate case flagged in Remark~\ref{rem:exactfail}. That's why minimizing $\gamma$ directly only gives a subgradient, while fixing $\gamma$ slightly above optimal (next paragraph) keeps the problem non-degenerate and restores an exact gradient. Might help to say this trade-off explicitly in the main text, since it's the reason the paper introduces the fixed-$\gamma$ SDP at all.}

\emph{Fixed-$\gamma$ SDP: exact gradients.} Fix the attenuation $\gamma$
strictly above the optimal level (the problem \eqref{eq:Hinf_cont}) can be solved for $\gamma^\star$, then $\gamma$ is picked larger with some gap). Rather than minimizing $\mu$, one then
minimizes the worst-case cost itself,
$\operatorname{tr}(XW_\theta)=\operatorname{tr}(Y^{-1}W_\theta)$, over the BRL LMI at
that fixed $\gamma$; with an epigraph slack $T$ and a Schur complement this is the SDP
\begin{equation}\label{eq:Hinf_fixed}
\begin{aligned}
V_{\mathrm{wc}}(\theta, \gamma)=&\min_{Y,L,T}\ \operatorname{tr}(T)\\
\text{s.t. } & Y\succ0,\quad
\begin{bmatrix} T & E_\theta\\ E_\theta & Y\end{bmatrix}\succeq0,\\
& \begin{bmatrix}
-Y & \star & \star & \star \\
0 & -\gamma^2 I & \star & \star \\
A_\theta Y + B_\theta L & E_\theta & -Y & \star \\
C_z Y + D_z L & 0 & 0 & -I
\end{bmatrix} \preceq 0,
\end{aligned}
\end{equation}
whose last LMI is exactly that of \eqref{eq:Hinf_cont} with $\mu=\gamma^2$ held
constant, and whose Schur slack enforces
$T\succeq W_\theta^{1/2}Y^{-1}W_\theta^{1/2}$, so $\operatorname{tr}(T)$ upper-bounds and at the optimum equals $\operatorname{tr}(Y^{-1}W_\theta)=\operatorname{tr}(XW_\theta)$.

\begin{corollary}[$H_\infty$ fixed-$\gamma$ design: exact gradient]\label{cor:hinf_fixed}
Fix $\gamma$ strictly above the optimal attenuation. Then \eqref{eq:Hinf_fixed} is
strictly feasible with a nondegenerate, strictly complementary dual, so
Assumption~\ref{as:env} holds and, by Theorem~\ref{prop:envelope} (equivalently
Proposition~\ref{prop:regexact}, the Schur slack being one added convex constraint),
$V_{\mathrm{wc}}(\theta, \gamma)$ is differentiable at $\theta$, with the design
gradient read from the BRL dual $S^\star_\gamma$ of \eqref{eq:Hinf_fixed},
\begin{equation}\label{eq:gradHinf_fixed}
\frac{\partial V_{\mathrm{wc}}}{\partial A_\theta}=2\,S_{\gamma,31}^\star\,Y^\star,\quad
\frac{\partial V_{\mathrm{wc}}}{\partial B_\theta}=2\,S_{\gamma,31}^\star\,(L^\star)^\top,
\end{equation}
of the same form as the minimal-$\gamma$ expression \eqref{eq:gradHinf} but now exact gradients rather than mere subgradients.
\end{corollary}

\emph{Estimation SDP (fixed-$\gamma_e$): exact gradient.} Exactly as
\eqref{eq:Jest} is the transpose dual of \eqref{eq:Jcont} under
$(A,B,Q,R,W)\to(A^\top,C^\top,W,V,Q)$, fixing $\gamma_e$ strictly above the
optimal filter attenuation and applying the same swap to \eqref{eq:Hinf_fixed}
gives the estimation analog. That is, we also swap $(E_\theta, C_z, D_z) \to (Q^{1/2}, [\,W_\theta^{1/2}\;\,0_{r_x \times r_u}\,]^\top, [0_{r_u \times r_x}\;\,V_\theta^{1/2}]^\top)$. The optimal filter gain $G^\star$ is the transpose of that returned ($K^\star$) in the control case. Similarly, the jacobians w.r.t. $A_\theta$ and $C_\theta$ are the transpose of the form of the jacobians in \eqref{eq:Hinf_fixed} w.r.t. $A_\theta$ and $B_\theta$, respectively.

\emph{Scalable game-Riccati route.} Exactly as the $H_2$ control SDP~\eqref{eq:Jcont}
has an $O(r_x^3)$ Riccati twin ($J_{\mathrm{cont}}^\star=\operatorname{tr}(PW)$ from the
control DARE), the fixed-$\gamma$ SDP~\eqref{eq:Hinf_fixed} admits a Riccati twin that
avoids the interior-point solve entirely: useful at scale, and a check on the SDP
dual near the minimal-$\gamma$ optimum, where the BRL LMI \eqref{eq:Hinf_cont}
loses rank (Corollary~\ref{cor:hinf}) and its dual is numerically delicate.

%=======================================================================
\section{The locally linearized case: MPC and the eKF}\label{sec:ekf}
%=======================================================================

The Bayesian filter is infinite dimensional; we approximate it by an extended Kalman
filter (eKF) \cite{anderson2012optimal,ramadan2024extended} propagating the first
two moments $x_{k\mid k},\Sigma_{k\mid k}$ (recursions in
Appendix~\ref{app:infostate}).

Concretely, with Jacobians $F_k = \partial f_\theta/\partial x|_{x_{k\mid k}}$ and
$H_k = \partial h_\theta/\partial x|_{x_{k\mid k-1}}$, the eKF alternates a
prediction and a measurement update steps similar to the Kalman filter.

If we stack the mean and the vectorized
upper-triangular elements of the state covariance into the information state vector
$\mathcal X_k = [x_{k\mid k}^\top,\, \operatorname{upvec}\Sigma_{k\mid k}^\top]^\top \in
\mathbb{R}^{n_\mathcal X}$, $n_\mathcal X = r_x + r_x(r_x+1)/2$, with lifted dynamics
induced by the (prediction-mode) eKF map (see Appendix~\ref{app:infostate} for more details)
\begin{equation} \label{eq:eKF as a dynamic system}
\mathcal{X}_{k+1}
     = \mathcal{F}_\theta(\mathcal{X}_k, u_k),
\end{equation}
and then linearize about the nominal information state (e.g. $(\mathcal{X}_0, u_0)$, since $k=0$ is current-time in receding-horizon convention)
\begin{equation*}
    \mathcal A_\theta = \left . \frac{\partial \mathcal F_\theta (\mathcal{X}, u_0)}{ \partial \mathcal{X}} \right |_{\mathcal{X} = \mathcal{X}_0},\quad \mathcal B_\theta = \left . \frac{\partial \mathcal F_\theta (\mathcal{X}_0, u)}{ \partial u} \right |_{u = u_0},
\end{equation*}
we end up in a linear system in the information state. The second-stage problem becomes $\widehat J_{sc} ^\star$, the approximate version of $J_{\mathrm{sc}}^\star$ in \eqref{eq:Jstoch_N},
\begin{equation}\label{eq:mpc}
\begin{aligned}
\widehat J^\star_{sc}(\theta) = &\min_{u_{0:N-1}}\
\frac{1}{N}\sum_{k=0}^{N-1} \mathcal L(\mathcal X_k,u_k) + \mathcal L^N(\mathcal X_N)\\
\text{s.t. } & \mathcal X_{k+1} = \mathcal A_\theta \mathcal X_k +
B_\theta u_k,\\
& \Sigma_{k\mid k} \succ 0,\quad \mathcal C(\mathcal X_{k+1},u_k) \le 0,
\end{aligned}
\end{equation}
where $\mathcal L$ now reconstructs the covariance term in $\operatorname{tr}(Q\Sigma_{k\mid k})$ from its vectorized upper-triangular values found in $\mathcal X_k$ (and similarily for the conic constraint $\Sigma_{k\mid k} \succ 0$). 
The function $\mathcal C(\cdot)$ represents the state/input constraints.

\emph{Design gradients from the dynamics costates.} Let $\lambda_k$ be the
multipliers (costates) of the dynamics equalities in \eqref{eq:mpc} in the corresponding Langrangian. By Theorem~\ref{prop:envelope}, the only explicit
$\theta$-dependence is through $\mathcal A_\theta, B_\theta$,
so
\begin{equation}\label{eq:gradmpc}
\frac{\partial \widehat J^\star_{sc}}{\partial \theta_i} = \sum_{k=0}^{N-1} \lambda_k^{\star\,\top} \Big(
\frac{\partial \mathcal A_\theta}{\partial \theta_i}\mathcal X_k^\star +
\frac{\partial \mathcal B_\theta}{\partial \theta_i} u_k^\star \Big).
\end{equation}
\suggestnote{Double-check the time index on $\lambda_k$ here against how the costate is actually defined in your Lagrangian: with the dynamics equality $\mathcal X_{k+1}=\mathcal A_\theta \mathcal X_k + \mathcal B_\theta u_k$ written as a constraint at "time $k$", the standard adjoint convention often attaches the multiplier as $\lambda_{k+1}$ (the multiplier introduced when eliminating $\mathcal X_{k+1}$), not $\lambda_k$. Not necessarily wrong -- conventions differ -- but worth a one-line definition of $\lambda_k$ right before this equation so readers can check the indexing themselves.}

\begin{rem}
    If the state evolution and the measurement operator $f_\theta$ and $h_\theta$ represent a linear system, then $\mathcal{A}_\theta = \operatorname{blkdiag}\,(A_\theta, \bar A _\theta)$ (the evolution of the error covariance is independent of the state estimate) and $\mathcal{B}_\theta = [B_\theta^\top \, 0]^\top$ (the error covariance is independent of the input). This is a consequence of the separation principle between estimation and control for fixed $\theta$.
\end{rem}

\begin{rem}[The structure of the controller as a dual control]
    The MPC control \eqref{eq:mpc} is feedback in the information state $\mathcal{X}_k$, instead of $x_{k \mid k}$ only as in the SDP case. This is due to the fact that the separation principle does not hold in the nonlinear regime: the input and state estimate can have an effect on the evolution of the estimation error covariance. Hence, the resulting controller admits probing/caution effects, the characteristics of a dual control. It is possible to turn \eqref{eq:mpc} into a function of $x_{k \mid k}$ only by fixing the value of $\Sigma_{k \mid k}$ to, for example, $\lambda_{\mathrm{probing}} I$, turning $\lambda_{\mathrm{probing}} \geq 0$ into a probing (excitation for data collection) knob an operator can choose based on the circumstances.
\end{rem}

\begin{rem}[Where exactness stops: the nonlinear case]\label{rem:ekfbreaks}
Theorem~\ref{prop:envelope} applies verbatim to the SDPs, which are genuinely
convex. In the nonlinear regime it applies to the convex program \eqref{eq:mpc},
not to the underlying stochastic control value \eqref{eq:Jstoch_N}. Three gaps
separate the two: (i)~Filtering
approximation: the eKF is a finite-dimensional approximation to the exact Bayesian recursion
\eqref{eq:Tupdate} \eqref{eq:Mupdate} by its first two moments, exact only for
affine $f_\theta,h_\theta$ \cite{anderson2012optimal} (the Kalman filter case).
(ii)~the prediction-only mode of the eKF used along the MPC horizon, which drops the innovation feedback. (iii)~Further linearization: $\mathcal A_\theta,
\mathcal B_\theta$ are taken about the current-time ($k=0$) information trajectory, so
\eqref{eq:gradmpc}. Nevertheless, Section~\ref{sec:numerical} shows that in practice this approximate gradient can still be used effectively in design optimization.
\end{rem}

%=======================================================================
\begin{rem}[Sensor allocation]\label{rem:sensor}
Deciding which of the candidate sensors to deploy fits ContEst without
disturbing the validity of its gradients. (a)~Adding binary indicators $b\in\{0,1\}^{\#\text{sensors}}$ and
costs $c$ to the upper level, $\min_{\theta,b}
J_{\mathrm{sc}}^\star(\theta,b)+c^\top b$, leaves the inner problem, and hence the
envelope gradients \eqref{eq:chain}, \eqref{eq:gradmpc}, unchanged. The price is a
(NP-hard) mixed-integer program \cite{joshi2009sensor,tzoumas2016sensor}.
(b)~Relaxing $b$ to continuous gains $\alpha\ge0$ ($C_{\alpha,\theta}=\operatorname{diag}(\alpha)\,C_{\mathrm{base},\theta}$) with an $\ell_1$ penalty in the design cost,
\begin{equation}\label{eq:Jdes}
J_{\mathrm{des}}^\alpha(\theta) = J_{\mathrm{des}}(\theta) + \lambda_\alpha \| \alpha \|_1,
\quad \alpha \ge 0,
\end{equation}
promotes sparse sensing \cite{dhingra2014admm,saraf2017h2,das2017sparse}. (c)~The $\ell_1$ kink is removed
by its epigraph lift \cite{boyd2004convex}, $\lambda_\alpha \mathbf 1^\top t$ with
$\alpha\preceq t$, keeping $J$ differentiable so the BFGS guarantees of
Remark~\ref{rem:algorithm} hold. Finally, a threshold on $\alpha^\star$ recovers a binary
deployment.
\end{rem}

\begin{rem}[Outer optimization]\label{rem:algorithm}
Since the design gradient is available,
Problem~\ref{prob:bilevel} is solved by a first-order method, here a projected
BFGS iteration on $\theta$ over the box $\Theta$
\cite{nocedal2006numerical}. Because the gradients are
exact (not finite-difference), the
secant pairs driving the BFGS update are accurate. Under the usual smoothness
and local strong-convexity assumptions the outer iteration attains the local
superlinear rate of full BFGS \cite{nocedal2006numerical}. When the inner value
is only piecewise smooth (e.g.\ the $\ell_1$-selection or worst-case $H_\infty$
regimes) and the returned gradient is a subgradient, a nonsmooth BFGS variant
with a suitable line search still converges reliably in practice
\cite{lewis2013nonsmooth}.
\end{rem}



\section{On the computational complexity of ContEst with interior-point methods}
\label{sec:sim:complexity}
%-----------------------------------------------------------------------

Because the envelope theorem (Section~\ref{sec:envelope}) yields the design
gradient from the primal and dual of a single inner solve, never by
differentiating through the solver's iterations, the per-design gradient is one inner
solve plus a matrix multiplication, and the outer BFGS iteration inherits the
complexity class of the inner solver.

\emph{Interior-point covariance/LMI solves scale as $O(r_x^6)$.} Both ContEst
inner problems place a covariance matrix among the decision variables. In the LQG
SDPs of Section~\ref{sec:lqg} the variable is the $r_x\times r_x$ covariance
$\Sigma$ (and its dual), i.e.\ $m=O(r_x^2)$ scalar unknowns; a primal-dual
interior-point method assembles and factorizes a dense Schur-complement (Newton)
system of size $m\times m$, at cost $O(m^3)=O(r_x^6)$ per iteration, times the
$O(\sqrt{r_x})$/$O(1)$ iterations interior-point methods take in practice. Empirically, solving the LQR
covariance program with the Clarabel interior-point solver took
$\approx\!2.6\,$s at $r_x=20$ and $\approx\!30\,$s at $r_x=30$: an
$(30/20)^6\!\approx\!11\times$ growth, matching the $O(r_x^6)$ prediction, so a
naive SDP inner problem caps ContEst at a few tens of states. To scale beyond that: (i) model structure and sparsity patterns can be exploited to reduce the number of nonzero entries in $\Sigma$, (ii) Krylov methods can be used within interior point implementations, instead of using direct Cholesky factorizations \cite{gondzio2012interior}.

\emph{The eKF--MPC inner problem meets the same wall.} There the decision trajectory
is the information state $\zeta=[x;\operatorname{upvec}\Sigma]$ of dimension
$r_x+\tfrac{r_x(r_x+1)}{2}=O(r_x^2)$; the horizon-$N$ MPC \eqref{eq:mpc} thus has
$O(N r_x^2)$ variables and its interior-point solve again grows steeply in $r_x$
(we measured $\approx\!33\,$s per evaluation at $r_x=32$).
Consequently the nonlinear/constrained studies of
Sections~\ref{sec:sim:distill}--\ref{sec:sim:pll} are run at modest $r_x$, and
large $r_x$ similarly requires model reduction, structure exploitation, and Krylov methods.


%=======================================================================
\section{Simulation studies}\label{sec:numerical}
%=======================================================================
\suggestnote{This entire section (and the whole paper) reports results only in tables -- there isn't a single figure. For a numerics-heavy section like this, at least one plot would likely help a lot: e.g.\ (a) a bar or line plot of the $J_{\mathrm{est}}/J_{\mathrm{cont}}/J_{\mathrm{tot}}$ reductions across the four studies (Table~\ref{tab:formulations} as a figure), (b) closed-loop state/estimation-error trajectories before vs.\ after co-design for one study, or (c) the MTDC sensor-count frontier of Table~\ref{tab:mtdc-frontier} as an actual frontier plot, which would make the "buy back sensors" headline land much faster than reading the table. Not a correctness issue, just a strong presentation suggestion given how central these results are to the paper's contribution.}

We evaluate ContEst on four co-design problems spanning very different physics. These studies are then followed by a complexity analysis of the two inner problems
under interior-point solvers (Section~\ref{sec:sim:complexity}).
\errornote{Section ordering mismatch: this sentence says the case studies are "followed by" the complexity analysis, but Section~\ref{sec:sim:complexity} is physically placed \emph{before} this Simulation Studies section (Sec.~\ref{sec:sim:complexity} is Sec.~6, this is Sec.~7). Same issue recurs later, e.g.\ the MTDC study says "exercising the structure-exploitation strategy flagged only in the abstract in Section~\ref{sec:sim:complexity}" as if that section comes later. Either swap the order of the two sections so complexity follows the case studies (matches the prose), or reword these references to say "preceded by" / "discussed above in Section~6".} All these examples were solved by the Clarabel
interior-point solver. The outer loop is the projected quasi-Newton method of
Remark~\ref{rem:algorithm} (BFGS) over the box $\Theta$, run from $5$ initial
points (the baseline design $\theta_{\mathrm{nom}}$ plus four uniform random
restarts in $\Theta$) keeping the best minimizer. The envelope-theorem gradients are exact, so we do not verify them
per study, but as a one-time sanity check we confirmed that
the analytic gradient matches central finite differences to a relative error below
$10^{-4}$.
\suggestnote{This blanket claim ("we do not verify them per study") sits awkwardly next to several places later where you \emph{do} report per-study gradient checks: the PLL trigger-1 stationary point (verified to $10^{-5}$), the MTDC block-diagonal control dual (verified to $0.55\%$), and the MTDC banded restriction that you flag as failing the check (over $100\%$ disagreement, "including the sign"). Consider softening this sentence to something like "...so we do not re-verify them in every study by default, except where flagged below" so it doesn't read as contradicting the later per-study checks.}\footnote{ContEst obtains the exact design gradient at essentially the cost of a
single inner solve, while forward (central) differences need $r_\theta+1$ ($2r_\theta$)
inner solves for an $r_\theta$-vector $\theta$, for an approximate gradient.}

Code reproducing every study is publicly
available.\footnote{\url{https://github.com/anonymous/ContEst} (placeholder
repository; the final URL will be provided on acceptance).}

For the studies in which linearization takes place, we report Monte-Carlo estimates of the realized closed-loop cost to confirm each optimized
design improves the true cost, not merely the surrogate problem. Starting from independent initial states sampled from $p_0$ and
with sampled process/measurement noise,
\begin{equation}\label{eq:costsplit}
\begin{aligned}
J_{\mathrm{sc}} &= \frac{1}{M}\sum_{m=1}^M\sum_{k}^T\!\big(\|x_k\|_Q^2+\|u_k\|_R^2\big),\\
J_{\mathrm{est}} &= \frac{1}{M}\sum_{m=1}^M\sum_{k}^T\|x_k-\hat x_{k\mid k}\|_Q^2,
\end{aligned}
\end{equation}
with $M$ the number of Monte-Carlo sample trajectories and $T$ each trajectory's
horizon length. These three studies use $M=100$ sample
trajectories over a horizon of $T=30$ steps, while the receding-horizon
MPC itself plans over $N=12$ future steps.

%-----------------------------------------------------------------------
\subsection{Spacecraft attitude determination and control (ADCS)}
\label{sec:sim:adcs}
%-----------------------------------------------------------------------

Pointing accuracy on every space mission is ultimately set by both the attitude estimator (the sensor suite: star trackers, gyros)
and the reaction-wheel authority. On a spacecraft these compete for a shared,
tightly constrained power/mass budget: a heavier, higher-power star tracker
or gyro reports lower noise, and a larger wheel delivers more torque, but every
watt and kilogram spent on one is unavailable to the others. We pose this on the
steady-state SDP path of Sections~\ref{sec:lqg} and \ref{sec:hinf}: $H_2$
state-feedback control under a hard cap on the steady-state control-effort
covariance, paired with worst-case ($H_\infty$) robust estimation, the
actuation and sensing axes tied by the shared budget.

A rigid satellite with inertia matrix
$J=\operatorname{diag}(4,6,5)\,\si{kg.m^2}$ regulates attitude and rate to zero
about its pointing equilibrium.
The state $x=[\phi,\theta_a,\psi,\omega_x,\omega_y,\omega_z]$ ($r_x=6$) is attitude and body rates; the input $u$ is three reaction-wheel
torques ($r_u=3$); the outputs are a star tracker (attitude) and a rate gyro
($r_y=6$, $C=I_6$). The only nonlinearity is the gyroscopic coupling $\omega\times
(J\omega)$, whose Jacobian vanishes at the pointing equilibrium $\omega=0$;
a single linearization there yields three decoupled double integrators.
With $\Delta t=\SI{0.1}{s}$,
\begin{subequations}\label{eq:adcs-model} (the $+$ superscript denotes the next time-step in a recursion)
\begin{align}
(\phi,\theta_a,\psi)^{+} &= (\phi,\theta_a,\psi) + \Delta t\,(\omega_x,\omega_y,\omega_z),\\
\omega^{+} &= \omega + \Delta t\,J^{-1}\big(e_{\mathrm{rw}}\,u - \omega\times(J\omega)\big),\\
h_\theta &= \big[\phi,\theta_a,\psi, \omega_x,\omega_y,\omega_z\big]^\top,
\end{align}
\end{subequations}
with $V_\theta=\operatorname{diag}(v_0/\alpha_{\mathrm{st}}^2\,I_3,\,v_0/\alpha_{\mathrm{gyro}}^2\,I_3), \, v_0=0.20$
(better precision $\Rightarrow$ lower variance), disturbance $w_k\sim\mathcal N(0,W)$ with
$W=\operatorname{diag}(10^{-6}I_3,10^{-4}I_3)$, and cost weights
$Q=\operatorname{diag}(4I_3,I_3)$ (attitude weighted four times more heavily than
rate, so the star tracker is intrinsically more valuable than the gyro) and
$R=0.1\,I_3$ (actuation is bounded by a hard effort cap, below, rather than by $R$).
The three design parameters are reaction-wheel authority $e_{\mathrm{rw}}$, which
scales the input and so enters $B_\theta$, and the star-tracker and rate-gyro precisions
$\alpha_{\mathrm{st}},\alpha_{\mathrm{gyro}}$, which enter the measurement
covariance $V_\theta$.

The control side is priced by the $H_2$ covariance SDP~\eqref{eq:Jcont} (dense,
$r_x=6$) in the stationary state covariance $\Sigma$, the input-second-moment
surrogate $Z_0$, and $L=K\Sigma$: its value
$J_{\mathrm{cont}}^\star(\theta)=\operatorname{tr}(Q\Sigma)+\operatorname{tr}(RZ_0)$
equals the LQG cost, with the design gradient in $e_{\mathrm{rw}}$ read as an exact
gradient from the LMI dual~\eqref{eq:gradAB} (no solver differentiation). We add one
\emph{hard} constraint, $\operatorname{tr}(Z_0)\le u_{\max}$: since
$Z_0\succeq K\Sigma K^\top$ is exactly the stationary control-effort covariance (the
steady-state input second moment), this caps the steady-state control-effort
covariance directly (a linear LMI that an LQR/Riccati weighting cannot
express, because $R$ can only price effort, not bound it). We set
$u_{\max}=0.6\,\operatorname{tr}(Z_0^{\mathrm{unc}})$ (60\% of the unconstrained
baseline effort). The sensing side is priced by the fixed-$\gamma$
$H_\infty$ robust-filter SDP at attenuation
$\gamma^2=4$, chosen as five times the minimal feasible $\gamma^2_{\min}=0.8$ so the
BRL LMI is feasible and its dual nondegenerate
(Corollary~\ref{cor:hinf_fixed}). Because $e_{\mathrm{rw}}$ enters only $B_\theta$ and the precisions only $V_\theta$, the two inner
values are coupled solely through the shared budget in the design cost (a
design-level) rather than dynamics-level coupling.

The design cost is a shared budget $c_b(e_{\mathrm{rw}}+\alpha_{\mathrm{st}}+
\alpha_{\mathrm{gyro}}-B)^2$ ($c_b=5$, $B=3$), with box
$e_{\mathrm{rw}}\in[0.3,3]$, $\alpha_{\mathrm{st}},\alpha_{\mathrm{gyro}}\in[0.3,5]$,
and baseline $\theta_{\mathrm{nom}}=\mathbf 1$ (budget spent uniformly). On the SDP
path actuation is bounded by the hard control-effort-covariance cap
$\operatorname{tr}(Z_0)\le u_{\max}$ rather than by pointwise input constraints on
the trajectory.

All five starts converged to the same $\theta$.
The co-design (Table~\ref{tab:adcs}) reallocates the budget:
it upgrades the star tracker most ($\alpha_{\mathrm{st}}\!:\,1\to1.70$), holds wheel
authority essentially fixed ($e_{\mathrm{rw}}\!:\,1\to1.005$ (the hard effort cap
saturates the return on actuation, so extra authority buys nothing), and drives the lightly weighted
rate gyro to its floor ($\alpha_{\mathrm{gyro}}\!:\,1\to0.30$), spending that budget
where it pays most. Improvements in control, estimation and design costs are reported in Table~\ref{tab:formulations}.

\begin{table}[t]
\centering
\caption{Spacecraft ADCS ($H_2$ control under a control-effort-covariance cap $+$
$H_\infty$ robust-estimation SDP, single
linearization about the pointing equilibrium): baseline (uniform budget) vs.\
co-designed allocation of the shared power/mass budget
$e_{\mathrm{rw}}+\alpha_{\mathrm{st}}+\alpha_{\mathrm{gyro}}\approx B=3$.}
\label{tab:adcs}
\begin{tabular}{llcc}
\hline
parameter & axis & baseline & optimal \\
\hline
$e_{\mathrm{rw}}$       & $f$ (reaction wheels) & 1.000 & 1.005 \\
$\alpha_{\mathrm{st}}$  & $V$ (star tracker)    & 1.000 & 1.703 \\
$\alpha_{\mathrm{gyro}}$& $V$ (rate gyro)       & 1.000 & 0.300 \\
\hline
\end{tabular}
\end{table}

%-----------------------------------------------------------------------
\subsection{Binary distillation column}
\label{sec:sim:distill}
%-----------------------------------------------------------------------

Distillation is the workhorse separation of the process
industries and its control is dominated by two classical layout decisions: on
which tray to introduce the feed, and on which tray to place the
inferential (temperature) sensor that is used to recover the composition
profile. These locations shape $f_\theta$ and $h_\theta$, so this is a ContEst problem.

A twelve-stage column (stage~1 reboiler
until stage~12 condenser) separates a light/heavy mixture; the state
$x=[\Delta x_1,\dots,\Delta x_{12}]$ ($r_x=12$) collects stage light-component
mole-fraction deviations about the nominal profile, in deviation form so
$f_\theta(0,0)=0$ for every $\theta$. The nonlinearity enters the dynamics twice: the vapour-liquid equilibrium
$y(x)=\alpha x/(1+(\alpha-1)x)$ ($\alpha=2.5$) and the bubble-point temperature
$T(x)=T_{B0}-\Delta T_B\,y(x)$. The single input is the reboiler boil-up
deviation $u=\Delta V$ ($r_u=1$), limited to $|\Delta V|\le1.5$; the measurements
($r_y=3$) are three Gaussian-weighted stage temperatures read at the sensor
locations. Integration is at $\SI{0.05}{s}$; the receding-horizon MPC plans over
$N=12$ steps. The design carries
$N_f=3$ feed-stage locations $\theta_f=(f_1,f_2,f_3)$ and $N_s=3$
temperature-sensor locations $\theta_h=(s_1,s_2,s_3)$, all in the interval $[2,11]$. Placement is nearly free to relocate,
so $J_{\mathrm{des}}$ carries only a mild regularisation
$\lambda_{\mathrm{des}}\lVert\theta-\theta_{\mathrm{nom}}\rVert^2$ about the
default $\theta_{\mathrm{nom}}$ with feeds at $(3,4.5,6)$ and sensors at
$(2,3.5,5)$.

With equilibrium
deviations $\Delta y_i=y(x_i^{\mathrm{ss}}+\Delta x_i)-y_i^{\mathrm{ss}}$, for
$i=1,\dots,12$ and $j=1,\dots,3$,
\begin{subequations}\label{eq:distill-model}
\begin{align}
\Delta x_i^{+} &= \Delta x_i + \tfrac{\Delta t}{M}\big(\mathrm{liq}_i+\mathrm{vap}_i\big),\\
\mathrm{liq}_i &= L_i(\theta_f)\,(\Delta x_{i+1}-\Delta x_i),\\
\mathrm{vap}_i &= (V+u)\,(\Delta y_{i-1}-\Delta y_i) + u\,(y_{i-1}^{\mathrm{ss}}-y_i^{\mathrm{ss}}),\\
h_{\theta,j} &= \textstyle\sum_{i=1}^{12}\hat q_i(s_j)\,\big(T(x_i^{\mathrm{ss}}+\Delta x_i)-T_i^{\mathrm{ss}}\big),
\end{align}
\end{subequations}
where the superscript $(\cdot)^{\mathrm{ss}}$ denotes the nominal steady-state
operating point about which the model is written in deviation form. The
boundary terms $\Delta x_{13}=\Delta y_0=y_0^{\mathrm{ss}}=0$, feed weights
$q_l(\theta_f)\propto \sum_{j=1}^{3}e^{-(l-f_j)^2/2\sigma_f^2}$ (normalized)
setting the internal liquid traffic
$L_i(\theta_f)=L_{\mathrm{ref}}+F\sum_{l>i}q_l$, and sensor weights
$\hat q_i(s_j)\propto e^{-(i-s_j)^2/2\sigma_{\mathrm{sens}}^2}$. Constants $\alpha=2.5$,
$L_{\mathrm{ref}}=2$, $V=2.5$, $F=2.5$, $M=1$,
$\sigma_f=\sigma_{\mathrm{sens}}=1.2$, $T_{B0}=\SI{380}{K}$,
$\Delta T_B=\SI{40}{K}$. The noise covariances $W=2\!\times\!10^{-5}I_{12}$,
$V=0.6\,I_3$ ($\si{K^2}$), and prior $\Sigma_0=10^{-2}I_{12}$. The cost weights
$Q_{ii}=20+30\,e^{-(i-6.5)^2/2\sigma_Q^2}$ ($\sigma_Q=2.4$, heaviest
mid-column), $R=2$,
$|u|_\infty\le1.5$, with a mild placement regularization
$\lambda_{\mathrm{des}}\lVert\theta-\theta_{\mathrm{nom}}\rVert^2$.

The design landscape is multimodal: the five-start BFGS converge to three distinct
minima, and the baseline start alone converges to an inferior one, so the random
restarts are what recover the reported optimum. ContEst spreads the feed
streams: two settle lower toward the reboiler ($f_{1,2}:3.0,4.5\to2.5,3.6$) while
the third climbs high up the column ($f_3:6.0\to9.0$), and raises the three
temperature sensors out of the bottom region across the mid-and-upper trays
($s:(2.0,3.5,5.0)\to(3.0,5.4,8.0)$), covering the composition profile rather than
clustering low (Table~\ref{tab:distill}). Improvements on the control, estimation and design costs are shown in Table~\ref{tab:formulations}.

\begin{table}[t]
\centering
\caption{Binary distillation: naive baseline vs.\ co-designed placements. All six
parameters are continuous tray locations (not gains); baseline $\theta_{\mathrm{nom}}$
places feeds at $(3,4.5,6)$ and sensors at $(2,3.5,5)$.}
\label{tab:distill}
\begin{tabular}{llcc}
\hline
parameter & axis & baseline & optimal \\
\hline
$f_1$ & $f$ (feed stage)    & 3.000 & 2.479 \\
$f_2$ & $f$ (feed stage)    & 4.500 & 3.584 \\
$f_3$ & $f$ (feed stage)    & 6.000 & 9.018 \\
$s_1$ & $h$ (temp.\ sensor) & 2.000 & 2.972 \\
$s_2$ & $h$ (temp.\ sensor) & 3.500 & 5.428 \\
$s_3$ & $h$ (temp.\ sensor) & 5.000 & 8.010 \\
\hline
\end{tabular}
\end{table}

%-----------------------------------------------------------------------
\subsection{Estimator co-design: PLL/DSE tuning in a low-inertia grid}
\label{sec:sim:pll}
%-----------------------------------------------------------------------

In this example we run ContEst as an online adaptive tool to co-design (tune) the state estimator itself: the per-inverter
dynamic state estimation (DSE) and the damping authority. Every
grid-following (GFL) inverter estimates grid frequency and phase with a
phase-locked loop (PLL), whose bandwidth is a classic bias-variance knob: a
narrow-band PLL is smooth but sluggish and lags real frequency excursions, while
a wide-band PLL tracks fast but passes more measurement noise into the frequency
estimate. In low-inertia, inverter-dominated grids, this tuning is decisive
for stability and for fast frequency response \cite{zhao2019power}. The PLL
bandwidth is therefore an estimator design variable that classical control
co-design cannot represent and sensor placement cannot capture.

Three GFL-inverter buses form a low-inertia chain. The state
$x=[\Delta\delta_{1:3},\Delta\omega_{1:3},\Delta\hat\omega^{\mathrm p}_{1:3}]$
($r_x=9$) stacks bus angle deviations, frequency deviations, and the per-bus PLL frequency estimate $\Delta\hat\omega^{\mathrm p}_i$; the input
$u=[\Delta P_{1:3}]$ is the inverter damping power ($r_u=3$). Each bus is measured
by a phasor measurement unit (PMU) angle channel and its PLL frequency channel
($r_y=6$). The genuine
nonlinearity is the synchronizing power $\propto\sin(\delta^{\mathrm{eq}}_{ij}+
\Delta\delta_i-\Delta\delta_j)$, in deviation form so $f_\theta(0,0)=0$ for
all $\theta$; integration is at $\SI{0.05}{s}$, the MPC plans over $N=12$ steps and
the closed-loop cost is Monte-Carlo sampled over $T=30$ steps (Eq.~\eqref{eq:costsplit}, $M=100$). The design parameters
are the per-bus inverter damping effectiveness
$\theta_f=(e_1,e_2,e_3)\in[0.3,3]^3$ and the per-bus PLL bandwidths
$\theta_h=(b_1,b_2,b_3)\in[2,20]\,\si{rad/s}$. 


Since the PLL bandwidths and damping gains are software values, and the useful allocation is operating-point dependent (the localized disturbance
fixes which bus is fastest and least certain), the low per-solve cost of
ContEst makes it natural to be used as an event-triggered adaptation layer that automatically tunes system parameters after sudden changes in the environment, such as the following scenario.

We exercise this adaptive layer over two triggers. First, under normal,
undisturbed operation (equilibrium, no post-fault kick, and a uniform
quiet-level prior $\Sigma_0$ at every bus), ContEst is run once from the naive
default $\theta_{\mathrm{nom}}$ ($e_i=1$, uniform $b_i=6$), producing a
\emph{peacetime} design $\theta_{\mathrm{peace}}$. Then a sudden event
concentrates a disturbance at bus~1 (large post-fault kick and an elevated
prior $\Sigma_0$ there), while buses~2,3 stay `quiet'; this second trigger
re-solves the \emph{same} co-design problem, now started from
$\theta_{\mathrm{peace}}$ rather than from $\theta_{\mathrm{nom}}$, producing
the post-fault design $\theta^\star$. The automatic design/tuning
must trade tracking speed against readout noise per bus in both triggers.

For $i=1,2,3$,
\begin{subequations}\label{eq:pll-model}
\begin{align}
&\Delta\delta_i^{+} = \Delta\delta_i + \Delta t\,\Delta\omega_i,\\
&\Delta\omega_i^{+} = \Delta\omega_i + \tfrac{\Delta t}{2H_i}
   \big(\!-P_i^{\mathrm e}(\Delta\delta) - D_i\Delta\omega_i + e_i u_i\big),\\
&\Delta\hat\omega^{\mathrm p,+}_i = \Delta\hat\omega^{\mathrm p}_i
   + \Delta t\,b_i\,(\Delta\omega_i - \Delta\hat\omega^{\mathrm p}_i),\\
&h_\theta = \big[\,\Delta\delta_j;\ \Delta\hat\omega^{\mathrm p}_j\,\big]_{j=1}^{3},
\end{align}
\end{subequations}
where $P_i^{\mathrm e}(\Delta\delta)=\sum_{j\neq i}K_{ij}\big[\sin(\Delta^{\mathrm{eq}}_{ij}
+\Delta\delta_i-\Delta\delta_j)-\sin\Delta^{\mathrm{eq}}_{ij}\big]$, with
$\Delta^{\mathrm{eq}}_{ij}=\delta^{\mathrm{eq}}_i-\delta^{\mathrm{eq}}_j$ the
equilibrium angle difference, is the synchronizing
power. Here $H=(1.6,2.2,2.0)\,\si{s}$ are the (low) inertia constants,
$D=(0.4,0.6,0.5)$ the natural dampings, $\delta^{\mathrm{eq}}=(0.15,-0.05,0.10)\,
\si{rad}$ the equilibrium angles, and the line synchronizing coefficients are
$K_{12}=1.2$, $K_{23}=1.1$, $K_{13}=0.4$ (an open $1$--$2$--$3$ chain with a weak
$1$--$3$ tie). The third line is the first-order PLL tracking the true
frequency with bandwidth $b_i$. The process noise covariance $W=\operatorname{blkdiag}
(2\!\times\!10^{-5}I_3,5\!\times\!10^{-4}I_3,5\!\times\!10^{-4}I_3)$ and the prior
$\Sigma_0$ (diagonal) large entry (1,1) at bus~1 ($3\!\times\!10^{-2}$), small elsewhere
($6\!\times\!10^{-3}$). The cost weights $Q=\operatorname{blkdiag}(15I_3,120I_3,0_3)$
(weighting the true angle and, heavily, the true frequency; PLL states
unweighted), $R=I_3$, $|u|_\infty\le1.5$. The design cost 
$J_{\mathrm{des}}=\lambda_e\sum_i(e_i-1)^2+\lambda_b\sum_i(b_i-b_{\mathrm{nom}})^2$,
$\lambda_e=0.5$, $\lambda_b=0.02$, $b_{\mathrm{nom}}=6$. The design coupling is
that the PLL bandwidth $b_i$ enters \emph{both} $f_\theta$ (the tracking line above)
and the measurement covariance, which is design-dependent,
\begin{equation}\label{eq:pll-V}
V(\theta)=\operatorname{blkdiag}\!\big(\sigma_\delta^2 I_3,\
\operatorname{diag}_i\,\sigma_\omega^2(1+\kappa\,b_i)\big),
\end{equation}
$\sigma_\delta^2=4\!\times\!10^{-3}$, $\sigma_\omega^2=2\!\times\!10^{-3}$,
$\kappa=0.30$: a wider-band PLL tracks faster but yields a noisier frequency
readout---the bias--variance knob $\theta_h$ resolves.

\emph{Trigger 1 (peacetime).} Under normal operation all five starts converge
to the same minimizer $\theta_{\mathrm{peace}}$. With little disturbance for
the actuator to counteract, the inverter damping is left exactly at its
baseline ($e_i=1$)---a genuine interior stationary point of the design cost,
confirmed by both the analytic and central-finite-difference gradients
vanishing there to within $10^{-5}$, not an optimizer stall---while every PLL
bandwidth is raised uniformly to trim ambient estimation error
($b_i\!:\,6\to10.6$). This cuts $J_{\mathrm{est}}$ by $12.1\%$
($38.0\to33.4$) and $J_{\mathrm{cont}}$ by only $4.2\%$ ($83.6\to80.1$), for a
$5.6\%$ reduction in the total $J_{\mathrm{tot}}$ ($121.5\to114.7$)---a modest,
almost purely estimation-side gain, exactly what a fault-free adaptation layer
should find: little reason yet to invest in actuation.

\emph{Trigger 2 (post-fault).} The bus-1 fault re-triggers ContEst, now
started from $\theta_{\mathrm{peace}}$ rather than from $\theta_{\mathrm{nom}}$
(all five starts, including one seeded at $\theta_{\mathrm{peace}}$ itself,
again converge to the same $\theta^\star$). ContEst reallocates sharply: it
raises the disturbed bus's PLL bandwidth further ($b_1\!:\,10.6\to13.7$) while
leaving the quiet buses' bandwidths essentially where trigger 1 left them
($b_{2,3}\!:\,10.6\to10.6$), and---now that there is a real disturbance to
counteract---unlocks the inverter damping authority for the first time
($e_1\!:\,1\to3.0$, $e_{2,3}\!:\,1\to1.95,2.19$; none of the bandwidths reach
the upper bound of $20$, because beyond the optimum the added readout noise in
$V(\theta)$ outweighs faster tracking; Table~\ref{tab:pll}). Measured against
$\theta_{\mathrm{peace}}$---the design actually in effect when the fault
hits, not the original naive $\theta_{\mathrm{nom}}$---this second trigger
cuts $J_{\mathrm{est}}$ by a further $3.0\%$ ($39.6\to38.4$) and
$J_{\mathrm{cont}}$ by $38.6\%$ ($223.2\to137.1$), for a further $31.6\%$
reduction in the total $J_{\mathrm{tot}}$ ($264.0\to180.7$). Measured instead
against the original naive baseline $\theta_{\mathrm{nom}}$---i.e.\ the
reduction if trigger 1 had never run---the combined effect of both triggers is
the $J_{\mathrm{est}}$ $-17.6\%$, $J_{\mathrm{cont}}$ $-39.5\%$,
$J_{\mathrm{tot}}$ $-33.9\%$ ($273.5\to180.7$) reported in
Table~\ref{tab:formulations}. Replacing the information-state MPC by
a certainty-equivalence controller (planning on the eKF mean as if it were exact)
yields an essentially identical realized cost ($J_{\mathrm{tot}}:273.3\to180.6$),
so on this near-separable study the improvement comes from the co-designed
parameters $\theta$ rather than from the covariance-aware controller structure.

Note that the optimal bandwidths sit strictly inside the box, which is
the signature of the bias-variance trade-off the framework now resolves
automatically, allocating the most estimator bandwidth exactly where the dynamics
are fastest and least certain.


\begin{table}[t]
\centering
\caption{Low-inertia grid (PLL/DSE): naive baseline $\theta_{\mathrm{nom}}$,
trigger-1 peacetime design $\theta_{\mathrm{peace}}$ (ContEst run under normal
operation), and trigger-2 post-fault design $\theta^\star$ (ContEst re-run,
started from $\theta_{\mathrm{peace}}$, once the bus-1 fault hits). The PLL
bandwidths $b_i$ enter both $f_\theta$ and the measurement covariance
$V(\theta)$.}
\label{tab:pll}
\begin{tabular}{llccc}
\hline
parameter & axis & $\theta_{\mathrm{nom}}$ & $\theta_{\mathrm{peace}}$ & $\theta^\star$ \\
\hline
$e_1$ & $f$ (inverter damping) & 1.000 & 1.000 & 3.000 \\
$e_2$ & $f$ (inverter damping) & 1.000 & 1.000 & 1.951 \\
$e_3$ & $f$ (inverter damping) & 1.000 & 1.000 & 2.191 \\
$b_1$ & $h,V$ (PLL bandwidth)  & 6.000 & 10.578 & 13.718 \\
$b_2$ & $h,V$ (PLL bandwidth)  & 6.000 & 10.569 & 10.567 \\
$b_3$ & $h,V$ (PLL bandwidth)  & 6.000 & 10.575 & 10.575 \\
\hline
\end{tabular}
\end{table}

%-----------------------------------------------------------------------
\subsection{Multi-terminal HVDC: droop and sparse voltage-sensor co-design}
\label{sec:sim:mtdc}
%-----------------------------------------------------------------------

Multi-terminal HVDC (MTDC) grids regulate DC-bus voltage with \emph{droop}
control: each converter reacts to its own local voltage deviation by
modulating its power injection, a fixed-structure, decentralized, proportional
law---the DC-grid analogue of frequency droop---rather than a freshly
synthesized feedback gain. Droop needs no communication, but an aggressive
droop setting on a meshed grid can excite a lightly-damped
\emph{inter-terminal resonance}, a worst-case (not average-case) failure mode
that $H_\infty$ prices directly. On the sensing side, not every terminal needs
a telemetered voltage sensor: the ring couples neighbouring buses, so an
unsensed terminal's converter current can be reconstructed from its
neighbours' voltages, trading estimation quality against a hard per-sensor
telemetry cost. This is the same $H_2$/$H_\infty$ SDP machinery of
Sections~\ref{sec:lqg}--\ref{sec:hinf}, at network scale ($r_x=30$) and under
an explicit sparsity pattern that reflects the physical topology---exercising
the structure-exploitation strategy flagged only in the abstract in
Section~\ref{sec:sim:complexity}. This study stays entirely on the SDP path:
no Riccati equation is solved or invoked anywhere below.

A ring of $N_T=15$ terminals has state $x=[\Delta V_i,\Delta I_i]_{i=1}^{N_T}$
($r_x=30$, bus-voltage and converter-current deviations), input $u$ the
per-terminal supplementary power-reference modulation ($r_u=15$), and output
$y$ the (possibly gain-scaled or absent) telemetered bus voltages
($r_y\le15$). With $\Delta t=\SI{5}{ms}$ and terminal $i$'s ring neighbours
$(i^-,i^+)$,
\begin{subequations}\label{eq:mtdc-model}
\begin{align}
\Delta V_i^+ &= \Delta V_i + \tfrac{\Delta t}{C_i}\big[G_{\mathrm{dc}}(\Delta V_{i^-}\!+\!\Delta V_{i^+}\!-\!2\Delta V_i) + \Delta I_i\big] + w_i,\\
\Delta I_i^+ &= \Delta I_i + \Delta t\,\big[u_i - \Delta I_i - k_0\theta_i\,\Delta V_i\big]/\tau_i,\\
h_{\theta,i} &= \alpha_i\,\Delta V_i,
\end{align}
\end{subequations}
where $\theta_i=k_i$ is the droop gain at terminal $i$ (entering $A$, the
$f_\theta$ axis) and $\alpha_i\in[0,1]$ is a per-terminal sensor gain
(entering $C$, the $h_\theta$ axis; $\alpha_i=0$ removes the sensor). Line
conductance $G_{\mathrm{dc}}=5$, droop leverage $k_0=12$, and heterogeneous
per-terminal capacitances $C_i$, converter lags $\tau_i$, and infeed
intensities $\sigma_i$ (driving the process-noise variance on $\Delta V_i$) are
fixed physical constants. Cost weights penalize voltage deviation far above
current ($Q=\operatorname{diag}(60,0.2,\dots)$, alternating per terminal),
$R=2I_{15}$, and a small per-sensor measurement variance $v=10^{-3}$.

\emph{Control: block-diagonal $H_\infty$.} $\theta$ enters $A$, so its
worst-case cost is priced by the fixed-$\gamma$ BRL SDP~\eqref{eq:Hinf_fixed}
at $\gamma^2=20$. A dense $(Y,L)$ solve is the $O(r_x^6)$ wall of
Section~\ref{sec:sim:complexity} in practice (1380 variables,
$\approx\!226\,$s); restricting $Y,L$ to a \emph{block-diagonal}
($2\times2$ per terminal, fully decentralized) pattern cuts this to 540
variables and $\approx\!0.4\,$s, a $+54.9\%$ conservatism cost in the reported
worst-case value but a clean, strictly complementary dual. A tighter
\emph{banded} restriction (each terminal's LMI block also coupling to its ring
neighbours, the pattern that would literally mirror the physical topology) is
numerically \emph{dual-degenerate} here---its envelope gradient disagrees with
a central finite difference by over $100\%$, including the sign, so it cannot
drive a gradient-based search---so the block-diagonal restriction is used for
control throughout, its dual verified against central differences on all 30
design parameters to $0.55\%$ relative error.

\emph{Sensing: banded $H_2$ with an $\ell_1$ sparsity penalty.} The estimation
SDP~\eqref{eq:Jest} has no such degeneracy: restricted to the \emph{banded}
(ring-neighbour) pattern that does mirror the physical topology, it matches
the full dense SDP's true cost $\operatorname{tr}(Q\Sigma^\varepsilon)$ to
within $+0.06\%$ at a fraction of the variables, and its dual gives the exact
gradient in both $\theta$ (through $A$) and the sensor gains $\alpha$ (through
$C$). An $\ell_1$ penalty $\lambda_s\lVert\alpha\rVert_1$ ($\lambda_s=0.05$)
in the design cost promotes sparse sensor placement, in the smooth epigraph
form of Section~\ref{sec:sim:adcs}'s reasoning (here $\alpha\ge0$ already, so
the penalty is simply linear).

Because $\theta$ enters both inner values through $A$, the two axes are
dynamically coupled; because $\alpha$ enters only $C$, the joint search
alternates a $\theta$-step (fixed $\alpha$, both SDPs) and an $\alpha$-step
(fixed $\theta$, estimation SDP only), each well-conditioned, rather than one
monolithic 30-variable solve. The design cost also prices droop stress,
$c_k\sum_i\theta_i$ ($c_k=5\times10^{-3}$), with box
$\theta_i\in[0.5,4]$, $\theta_{\mathrm{nom}}=3$ (uniform droop, all 15 sensors
on).

All five starts converged to a design within $1.5\%$ of each other; the best
retains 12 of 15 sensors (dropping terminals 2, 6, 8) and spreads droop over
$\theta^\star\in[0.89,1.67]$, well below the uniform baseline. This cuts the
estimation cost $J_{\mathrm{est}}$ by $31.8\%$ ($0.650\to0.443$) and the
control cost $J_{\mathrm{cont}}$ by $33.4\%$ ($2.476\to1.650$), for a $33.0\%$
reduction in $J_c=J_{\mathrm{cont}}+J_{\mathrm{est}}$ ($3.126\to2.093$) and,
\errornote{Numbers don't match Table~\ref{tab:mtdc-frontier} below: this "best" design uses 12 of 15 sensors, but that table only has rows for 15/9/6/4/3 sensors, and none of its co-designed-column values equal $J_{\mathrm{est}}=0.443$ or $J_{\mathrm{cont}}=1.650$ (closest are the 9-sensor row: $0.460$/$1.645$, and the 15-sensor row: $0.392$/$1.596$). Please double check: either the 12-sensor point needs its own row in Table~\ref{tab:mtdc-frontier}, or clarify in the text that the frontier table sweeps a different (fixed-count) experiment than the freely-optimized 12-sensor result reported here, so readers don't try to cross-reference the two and get confused.}
after the droop-stress and sensor-sparsity price, a $21.6\%$ reduction in the
total $J_{\mathrm{tot}}$ ($3.351\to2.626$). As in the other studies, these are
the exact steady-state SDP/Lyapunov values, not Monte-Carlo estimates: unlike
ADCS's single linearization or the eKF-MPC studies' receding-horizon surrogate,
this plant is linear by construction with no unmodeled dynamics, so there is
no surrogate-vs-truth gap for a Monte-Carlo check to reveal.

\begin{table}[t]
\centering
\caption{MTDC sensor-count frontier: fixed-droop (baseline $\theta_{\mathrm{nom}}=3$)
greedy sensor pruning vs.\ co-designing the droop at each retained-sensor
count (same sensor sets both ways, so the gap isolates the co-design gain).}
\label{tab:mtdc-frontier}
\small
\begin{tabular}{lcccc}
\hline
& \multicolumn{2}{c}{$J_{\mathrm{est}}$} & \multicolumn{2}{c}{$J_{\mathrm{cont}}$} \\
\#sensors & pruned & co-des. & pruned & co-des. \\
\hline
15 & 0.650 & 0.392 & 2.476 & 1.596 \\
9  & 0.770 & 0.460 & 2.476 & 1.645 \\
6  & 0.927 & 0.560 & 2.476 & 1.645 \\
4  & 1.082 & 0.653 & 2.476 & 1.644 \\
3  & 1.170 & 0.697 & 2.476 & 1.645 \\
\hline
\end{tabular}
\end{table}

Table~\ref{tab:mtdc-frontier} is the headline result: co-designing the droop
for a \emph{fixed} sensor budget beats naive pruning at every count, and by
enough to buy back sensors outright---the co-designed droop with only \textbf{6}
sensors ($J_{\mathrm{est}}=0.560$) already matches the fixed-droop baseline's
\textbf{15}-sensor value ($J_{\mathrm{est}}=0.650$), i.e.\ $60\%$ of the
voltage telemetry is removed with no loss in estimation quality, purely by
retuning the actuation-side droop. A sensor-selection method with no actuation
axis \cite{joshi2009sensor,tzoumas2016sensor} cannot pose this trade: droop
reshapes the \emph{dynamics} the estimator has to track, so a better-tuned
droop makes the remaining sensors' information go further, exactly the
interior, cross-axis trade that motivates ContEst.

%-----------------------------------------------------------------------
\subsection{Summary of improvements}
\label{sec:sim:summary}
%-----------------------------------------------------------------------

Table~\ref{tab:formulations} consolidates the four studies (dimensions, the design
split $\theta=(\theta_f,\theta_h)$, the inner value, and
the reductions); the absolute baseline$\to$optimum costs are reported compactly with
each study above. In three of the four studies---ADCS and the two eKF-MPC
studies---the reported
numbers are Monte-Carlo averages of the realised cost ($M=100$ samples, $T=30$-step
rollouts), never the steady-state SDP/gramian values used only to drive the design
search; the fourth (MTDC) reports the exact steady-state SDP/Lyapunov values instead,
since its plant is linear by construction with no surrogate-vs-truth gap to check
(Section~\ref{sec:sim:mtdc}). The spacecraft ADCS study is the linear case, solved on
the steady-state SDP path ($H_2$ control under a hard control-effort-covariance cap
$+$ worst-case $H_\infty$ robust estimation): its sensing precisions and wheel
authority compete for a shared power/mass budget, and the co-design drives most of
the reduction through robust sensing while control also falls, because the
realised control cost is measured on the true state and so partly inherits the
estimation-error reduction ($J_{\mathrm{est}}$ $-42.7\%$, $J_{\mathrm{cont}}$
$-18.3\%$) for a $34.0\%$ reduction in the total. The two
nonlinear studies (distillation, PLL/DSE) use the eKF-MPC second stage; the PLL/DSE
study additionally demonstrates ContEst as an event-triggered adaptation layer, run
once under normal operation and re-triggered by the fault (Section~\ref{sec:sim:pll}).
The MTDC study
returns to the linear SDP path at network scale, restricting both inner values to a
sparsity pattern reflecting the ring topology; its headline is cross-axis rather than
purely a cost reduction---co-designing the droop lets $60\%$ of the voltage sensors be
removed at no loss in estimation quality (Table~\ref{tab:mtdc-frontier})---alongside a
$21.6\%$ reduction in the total design cost. Across all
four studies the total cost falls by $9$--$34\%$, with $J_{\mathrm{est}}$ reduced in
every case, and each optimum is an interior trade rather than a parameter pinned at a
bound.
Every study compares against the default baseline $\theta_{\mathrm{nom}}$ of its
section. For the eKF-MPC studies the baseline and the co-design run the \emph{same}
information-state MPC controller and differ only in the design $\theta$, so the
reported reduction isolates the co-design (design-parameter) gain rather than any
controller upgrade.

A caveat on interpretation: the baseline parameter values $\theta_{\mathrm{nom}}$
above were chosen to be plausible rather than carefully hand-tuned, so the
reported reductions should not be read as ContEst outperforming a state-of-the-art
design. The contribution is not a competing tuning heuristic but the exact,
dual-based design gradient that lets a first-order method \emph{search} for the
joint optimum at all; the improvements simply confirm that this gradient points
in a useful direction.


\begin{table*}[t]
\centering
\caption{Consolidated view of the four studies}
\label{tab:formulations}
\begin{tabular}{lcclcccc}
\hline
study & $(r_x,r_u,r_y)$ & $\theta$: $n_f{+}n_h$ (enters) & inner value
& $J_{\mathrm{est}}\!\downarrow$ & $J_{\mathrm{cont}}\!\downarrow$ & $J_{\mathrm{tot}}\!\downarrow$\\
\hline
ADCS & $(6,3,6)$   & $1{+}2$: $f_\theta$, $V_\theta$ & $H_2^{\mathrm{cap}} + H_\infty$ & $42.7\%$ & $18.3\%$  & $\mathbf{34.0\%}$\\
Distillation      & $(12,1,3)$  & $3{+}3$: $f_\theta$, $h_\theta$ & eKF-MPC & $35.3\%$ & $5.6\%$  & $\mathbf{9.5\%}$\\
PLL/DSE    & $(9,3,6)$   & $3{+}3$: $f_\theta$, $(f_\theta,V_\theta)$ & eKF-MPC & $17.6\%$ & $39.5\%$ & $\mathbf{33.9\%}$\\
MTDC & $(30,15,\le\!15)$ & $15{+}15$: $f_\theta$, $h_\theta$ & $H_\infty^{\mathrm{bd}} + H_2^{\mathrm{band}}$ & $31.8\%$ & $33.4\%$ & $\mathbf{21.6\%}$\\
\hline
\end{tabular}
\end{table*}

%=======================================================================
\section{Conclusion}\label{sec:conclusion}
%=======================================================================

We introduced ContEst, a co-design framework that brings state and parameter
estimation into control co-design by optimizing system and sensing parameters against a stochastic (dual-control) objective
expressed by the information state. The inner value takes a convex form in
three regimes: a pair of semidefinite programs in the average-case LQG ($H_2$)
setting, their worst-case $H_\infty$ analog (a bounded-real-lemma semidefinite program); and a convex information-state MPC
in the locally linearized nonlinear regime. Sensor selection enters as a
mixed-integer problem with a sparse $\ell_1$ relaxation whose duals price the
essential sensors. In every case the design gradient is obtained exactly and cheaply
from the inner dual variables, the LMI duals or dynamics costates, by the envelope
theorem, sidestepping both implicit differentiation of the solution map and
differentiation of the KKT system.

Future work includes chance-constrained and scenario-based inner problems, Bayesian filters beyond the eKF, and decomposition of the outer problem
along multidisciplinary-design-optimization architectures for large interconnected
systems.
\suggestnote{The conclusion doesn't mention the approximation gaps you were careful to spell out in Remark~\ref{rem:ekfbreaks} (eKF vs.\ exact Bayesian filter, prediction-only mode dropping innovation feedback, linearization about the current-time information state) or the degeneracy caveat of Remark~\ref{rem:exactfail}/Corollary~\ref{cor:hinf}. A one- or two-sentence limitations paragraph before "Future work" would round out the conclusion and preempt a reviewer asking "so when does this NOT work exactly."}

\begin{ack}
The authors thank colleagues at Argonne National Laboratory for helpful
discussions.
\end{ack}

\appendix

%=======================================================================
\section{The information state and its LQG and eKF specializations}\label{app:infostate}
%=======================================================================

The second stage is a stochastic (dual) control problem
\cite{feldbaum1960dual,bar1974dual}. A causal law must be a function of the
available information $\mathcal{Z}_k$, not of the unobserved $x_k$; the
sufficient statistic is the filtering density $p(x_k \mid \mathcal{Z}_k;\theta)$,
governed by the Bayesian filter, which alternates a time update,
\begin{equation}\label{eq:Tupdate}
p(x_{k+1} \mid \mathcal{Z}_k, u_k) = \int p(x_{k+1} \mid x_k, u_k;\theta)\,
p(x_k \mid \mathcal{Z}_k)\, dx_k,
\end{equation}
and a measurement update,
\begin{equation}\label{eq:Mupdate}
p(x_{k+1} \mid \mathcal{Z}_{k+1}) \propto p(y_{k+1} \mid x_{k+1};\theta)\,
p(x_{k+1} \mid \mathcal{Z}_k, u_k).
\end{equation}
Both \eqref{eq:Tupdate}--\eqref{eq:Mupdate} involve the output
equation~\eqref{eq:ss-output} and hence the sensing parameters in $\theta$. The
standard smoothing identity makes this explicit in the cost.

\begin{lem}[Smoothing {\cite[Ch.~10]{resnick2019probability}}]\label{lem:smoothing}
For a probability space $(\Omega,\mathcal{A},\mathbf{P})$, a measurable
$\chi:\Omega\to\mathbb{R}$ with $\mathbb{E}_{\mathbf P}|\chi|<\infty$, and
$\sigma$-algebras $\mathcal{A}_0 \subset \mathcal{A}_1 \subset \mathcal{A}$,
$\mathbb{E}_{\mathbf P}\{\chi \mid \mathcal{A}_1\} =
\mathbb{E}_{\mathbf P}\{\mathbb{E}_{\mathbf P}\{\chi \mid \mathcal{A}_0\} \mid
\mathcal{A}_1\}$.
\end{lem}

\begin{lem}\label{lem:stagecost}
For a causal law $u_k = u_k(\mathcal{Z}_k, p_0)$,
\begin{equation*}
\mathbb{E}\,\ell(x_k,u_k) = \!\iint \ell(x_k,u_k)\,
p(x_k \mid \mathcal{Z}_k;\theta)\, p(\mathcal{Z}_k;\theta)\, dx_k\, d\mathcal{Z}_k.
\end{equation*}
\end{lem}
\begin{proof}
Conditioning on $\mathcal{Z}_k$ yields a $\sigma$-algebra contained in the one
underlying \eqref{eq:Jstoch_N}; apply Lemma~\ref{lem:smoothing}.
\end{proof}

Therefore, the term $\E \| x_k \| _Q^2$ can be expressed by
\begin{align*}
    \E \| x_k \| _Q^2 &= \E \left \{ \E \left \{  \| x_{k}\|_Q^2 \mid \mathcal{Z}_k \right \} \right \}\\
    &=\E \left \{  \operatorname{tr}(Q \Sigma_{k \mid k}) +  \| x_{k \mid k}\|_Q^2 \right \},
\end{align*}
\explainnote{Spelling out the step that produces the second line above (the real content; the first line is just the tower property, Lemma~\ref{lem:smoothing}, applied to $\chi=\|x_k\|_Q^2$). The work is in the inner conditional expectation $\E\{x_k^\top Qx_k\mid\mathcal Z_k\}$. Write $x_k=x_{k\mid k}+\varepsilon_k$ with $x_{k\mid k}=\E\{x_k\mid\mathcal Z_k\}$ and $\varepsilon_k=x_k-x_{k\mid k}$, so $\E\{\varepsilon_k\mid\mathcal Z_k\}=0$ by construction (that's what "conditional mean" means). Expanding the quadratic form,
\[
x_k^\top Qx_k = x_{k\mid k}^\top Qx_{k\mid k} + 2\,x_{k\mid k}^\top Q\varepsilon_k + \varepsilon_k^\top Q\varepsilon_k,
\]
then take $\E\{\cdot\mid\mathcal Z_k\}$ term by term: $x_{k\mid k}$ is $\mathcal Z_k$-measurable, so the first term passes through unchanged; the cross term vanishes because $\E\{\varepsilon_k\mid\mathcal Z_k\}=0$; and for the last term, $\varepsilon_k^\top Q\varepsilon_k=\operatorname{tr}(Q\varepsilon_k\varepsilon_k^\top)$ (a scalar equals its own trace, then cyclic property), so
\[
\E\{\varepsilon_k^\top Q\varepsilon_k\mid\mathcal Z_k\}=\operatorname{tr}\big(Q\,\E\{\varepsilon_k\varepsilon_k^\top\mid\mathcal Z_k\}\big)=\operatorname{tr}(Q\Sigma_{k\mid k}),
\]
using the definition of $\Sigma_{k\mid k}$ in~\eqref{eq:filterMeanCov} (trace and conditional expectation commute since both are linear). Altogether $\E\{x_k^\top Qx_k\mid\mathcal Z_k\}=\|x_{k\mid k}\|_Q^2+\operatorname{tr}(Q\Sigma_{k\mid k})$, exactly the align* above. This is the same conditional bias--variance decomposition used throughout estimation theory (mean-square value = variance + squared mean), here with zero "bias" cross-term because $x_{k\mid k}$ is defined as the conditional mean itself.}
where the first two conditional moments are
\begin{equation}
\begin{aligned}\label{eq:filterMeanCov}
     x_{k \mid k} &= \E \left \{ x_k \mid \mathcal Z_{k}\right \},\\
    \Sigma_{k \mid k} &= \E \left \{ [x_k-x_{k \mid k}][x_k-x_{k \mid k}]^\top \mid \mathcal Z_{k} \right \},
\end{aligned}
\end{equation}
are substituted above after writing $x_k = x_{k \mid k} + \varepsilon_k$, where $\varepsilon_k$ is the estimation error (of zero mean), and using the cyclic property of the trace operator.

Hence, the cost functions \eqref{eq:Jstoch_N} and \eqref{eq:Jstoch_infty} can be represented by 
\begin{align*}
     &J_{\mathrm{sc}} = \frac{1}{N} \cdot \E \Bigg [ \| x_{N\mid N} \| ^2_ Q  + \operatorname{tr}(Q \Sigma_{N \mid N})  \nonumber\\
    &+ \sum_{k=0}^{N-1} \left [  \| x_{k\mid k} \|^2_ Q  + \|u_k \|_R^2 + \operatorname{tr}(Q \Sigma_{k \mid k}) \right ] \Bigg ],
\end{align*}
or in the infinite-horizon case
\begin{align*}
     &J_{\mathrm{sc}} = \lim_{k \to \infty} \E  \left [  \| x_{k\mid k} \|^2_ Q  + \|\kappa(\mathcal{Z}_k) \|_R^2 + \operatorname{tr}(Q \Sigma_{k \mid k}) \right ].
\end{align*}

In this paper we have two regimes, linear and nonlinear. For the linear cases (Sections~\ref{sec:lqg} and \ref{sec:hinf}), we adopt the infinite horizon cost, with $u_k = \kappa(\mathcal{Z}_k) = K x_{k \mid k}$, that is, certainty equivalent state-feedback control (due to the separation principle: the state error covariance $\Sigma_{k \mid k}$ evolution is not dependent on $u_k$. After substituting the feedback control, and using the cyclic property of the trace again, the cost becomes that in \eqref{eq:linear stoch cost}, given that $\Sigma^\epsilon = \lim_{k \to \infty} \E \varepsilon_k \varepsilon_k^\top = \lim_{k \to \infty} \Sigma_{k \mid k}$, $\Sigma = \lim _{k \to \infty} \E x_{k \mid k} x_{k \mid k}^\top$, the estimation error covariance and the state covariance (about the origin), respectively.
\explainnote{This is the promised "full derivation" of Eq.~\eqref{eq:linear stoch cost} referenced back in Section~4, but it's compressed into one clause ("using the cyclic property of the trace again") without being shown. Spelled out: start from the general infinite-horizon formula just above, $J_{\mathrm{sc}}=\lim_k\E[\|x_{k\mid k}\|_Q^2+\|u_k\|_R^2+\operatorname{tr}(Q\Sigma_{k\mid k})]$. Substitute $u_k=Kx_{k\mid k}$ and apply the trace identity $a^\top Ma=\operatorname{tr}(Maa^\top)$ to both quadratic forms,
\[
\|x_{k\mid k}\|_Q^2=\operatorname{tr}(Qx_{k\mid k}x_{k\mid k}^\top),
\]
\[
\|Kx_{k\mid k}\|_R^2=\operatorname{tr}\big(RKx_{k\mid k}x_{k\mid k}^\top K^\top\big).
\]
Take $\E[\cdot]$ termwise; trace and expectation are both linear, and $K$ is deterministic so it factors out:
\[
\E\|x_{k\mid k}\|_Q^2=\operatorname{tr}\big(Q\,\E[x_{k\mid k}x_{k\mid k}^\top]\big),
\]
\[
\E\|Kx_{k\mid k}\|_R^2=\operatorname{tr}\big(RK\,\E[x_{k\mid k}x_{k\mid k}^\top]\,K^\top\big).
\]
As $k\to\infty$, $\Sigma:=\lim_k\E[x_{k\mid k}x_{k\mid k}^\top]$ (zero-mean since $x_0,w_k,v_k$ all are) turns these into exactly $\operatorname{tr}(Q\Sigma)$ and $\operatorname{tr}(RK\Sigma K^\top)$. For the last term, $\Sigma_{k\mid k}$ is the Kalman-filter error covariance, which for an LTI model is deterministic (depends only on the model, not on realized $y_k$'s), so $\E[\operatorname{tr}(Q\Sigma_{k\mid k})]=\operatorname{tr}(Q\Sigma_{k\mid k})$ trivially and $\lim_k\Sigma_{k\mid k}=\Sigma^\varepsilon$ is just its steady-state value. Summing the three pieces gives $J_{\mathrm{sc}}=\operatorname{tr}(Q\Sigma)+\operatorname{tr}(RK\Sigma K^\top)+\operatorname{tr}(Q\Sigma^\varepsilon)=J_{\mathrm{cont}}+J_{\mathrm{est}}$, i.e.\ Eq.~\eqref{eq:linear stoch cost}. It checks out -- it's just easy to gloss over on a first read since none of these intermediate steps are actually shown.}

For the nonlinear regime, we choose the eKF to track the evolution of the state first two moments. In addition to \eqref{eq:filterMeanCov}, let
\begin{align*}
     x_{k \mid k-1} &= \E \left \{ x_k \mid Z_{k-1}, u_{k-1} \right \}, \\
    \Sigma_{k \mid k-1} &= \E \left \{ [x_k-x_{k \mid k-1}][x_k-x_{k \mid k-1}]^\top \mid Z_{k-1},u_{k-1} \right \}.
\end{align*} 
At each time step, the eKF is the standard recursion
\add{\cite{anderson2012optimal,jazwinski2007stochastic}},
\add{with $W=\Cov(w_k)\succeq0$ and $V=\Cov(v_k)\succ0$ the process and
measurement noise covariances already introduced in
Section~\ref{sec:formulation} (possibly $\theta$-dependent, as noted there),}
\begin{align}
 x_{k+1 \mid k} = f_\theta( x_{k \mid k},u_k),\quad \label{eq:eKF_state}
\Sigma_{k+1 \mid k} = F_k \Sigma_{k \mid k} F_k^\top + W,
\end{align}
\errornote{$W$ and $V$ are never redefined locally -- they were introduced once, all the way back in Sec.~3 ("the disturbances $w_k,v_k$ are i.i.d. ... with covariances $W,V$"), and then not mentioned again until they resurface here in Eq.~\eqref{eq:eKF_state}/\eqref{eq:eKF_supportVariables}, roughly 1000+ lines and 6 sections later. Same pattern as the $Q$ gap flagged earlier at Eq.~(4): a symbol defined once far upstream and reused cold in the appendix. A one-clause reminder (added above) is an easy fix; worth checking the rest of the appendix for other symbols reused this way without a local reminder (e.g. $Q$ itself reappears in \eqref{eq:eKF_supportVariables}'s surrounding derivation too).}
where
\begin{equation}
\begin{aligned} \label{eq:eKF_supportVariables}
& x_{k\mid k}= x_{k \mid k-1} + \Omega_{k} \left [y_{k}-h_\theta(x_{k \mid k-1}, u_k)\right],\\
&\Sigma_{k \mid k} = \left [I - \Omega_{k} H_{k} \right]\Sigma_{k \mid k-1},\\
 &\Omega_{k}= \Sigma_{k \mid k-1} H_{k}^\top \left [ H_{k} \Sigma_{k \mid k-1} H_{k}^\top+V\right]^{-1}, \\
 &F_k = \left. \frac{\partial f_\theta(x,u_k)}{\partial x} \right | _{x_{k\mid k}},\quad H_{k} = \left. \frac{\partial h_\theta(x,u_k)}{\partial x} \right | _{x_{k\mid k-1}}.
\end{aligned}
\end{equation}
The recursion is initialized by $ x_{0 \mid -1}$, $\Sigma_{0 \mid -1}$, the moments of $p_0$.

Note that the evolution of the eKF in \eqref{eq:eKF_supportVariables} depends on $y_k$. As an approximation, we drop this dependence, so we can approximate the evolution of the eKF to the future deterministically, before future measurements are received. This approximation can be plausible in many applications, since the measurement correction term $[y_{k} - h_\theta(x_{k\mid k-1}, u_k)]$, in \eqref{eq:eKF_supportVariables}, can be close to independent from $\Omega_{k}$, almost of a zero-mean white noise sequence when the eKF is fine-tuned \cite[Sec.~8.2]{anderson2012optimal}, and when $f_\theta$ is not highly nonlinear in its first argument. This weakens the need for the expectation in the finite-horizon $J_{\mathrm{sc}}$, and leads to a measurement-free version of the eKF in \eqref{eq:eKF_state} and \eqref{eq:eKF_supportVariables}. Now we can define the deterministic version of the eKF, propagating the information state $(x_{k \mid k}, \Sigma_{k \mid k})$,
\begin{equation*}
\mathcal{X}_{k+1} = 
\begin{pmatrix}
    x_{k+1 \mid k+1} \\
    \operatorname{upvec}\Sigma_{k+1 \mid k+1}
\end{pmatrix}
     = \mathcal{F}_\theta(\mathcal{X}_k, u_k),
\end{equation*}
where $\mathcal{F}_\theta$ is constructed by \eqref{eq:eKF_state} and \eqref{eq:eKF_supportVariables} when the term $y_k - h_\theta(x_{k \mid k-1},u_k)$ is set to zero, and $\operatorname{upvec}$ is the upper-triangle elements vectorized (to reduce the state dimension using symmetry). 



%=======================================================================
\section{Proof of the LQG regularity Proposition}\label{app:lqgreg}
%=======================================================================


\explainnote{You will get in trouble with the assumption labels and the equation labels, you have two A2, A3 etc}
This appendix proves Proposition~\ref{prop:lqgreg} by verifying
Assumption~\ref{as:env}(A1)--(A4) for the control SDP~\eqref{eq:Jcont} and, in the
course of (A3), deriving the dual identity $S_{11}^\star=P$. The estimation
SDP~\eqref{eq:Jest} is its formal dual (swap $A_\theta\!\leftrightarrow
A_\theta^\top$, $B_\theta\!\leftrightarrow C_\theta^\top$, $W\!\leftrightarrow Q$,
$R\!\leftrightarrow V$; detectability for stabilizability); the same steps yield
$P^{\varepsilon\star}=\Sigma^\varepsilon$ and $S_{11}^{\prime\star}=
\Sigma^\varepsilon$. Write the program as $V(\theta)=\min_z f(z,\theta)$ s.t.\
$z\in\mathcal C(\theta)$, with $z=(\Sigma,L,Z_0)$, linear objective
$f=\operatorname{tr}(Q\Sigma)+\operatorname{tr}(RZ_0)$, and $\mathcal C(\theta)$ the
two inequalities $M_1=\left[\begin{smallmatrix}Z_0&L\\ L^\top&\Sigma
\end{smallmatrix}\right]\succeq0$ and $M_2=\left[\begin{smallmatrix}\Sigma-W&
A_\theta\Sigma+B_\theta L\\ \star&\Sigma\end{smallmatrix}\right]\succeq0$. Assume
$(A_\theta,B_\theta)$ stabilizable and $Q,R,W\succ0$.

\emph{(A1) Convexity, smooth data, strong duality.} $f$ is linear and
$\mathcal C(\theta)$ is defined by inequalities affine in $z$ (hence a convex set)
whose coefficients are affine in the model matrices $(A_\theta,B_\theta,Q,R,W)$, so
the data are jointly $C^1$ in $(z,\theta)$. A strictly feasible point exists: for any
stabilizing $K_0$ let $\Sigma_0\succ0$ solve $A_{\mathrm{cl},0}\Sigma_0
A_{\mathrm{cl},0}^\top-\Sigma_0+W\prec0$ ($A_{\mathrm{cl},0}=A_\theta+B_\theta K_0$
Schur, $W\succ0$) and set $L_0=K_0\Sigma_0$, $Z_0\succ K_0\Sigma_0K_0^\top$; both
inequalities hold strictly. Slater then gives strong duality with attained primal
and dual solutions \cite{boyd2004convex}.

\emph{(A2) Unique primal minimizer.} With $W\succ0$, $M_2\succeq0$ forces
$\Sigma\succ0$ and (Schur complement) $A_{\mathrm{cl}}\Sigma A_{\mathrm{cl}}^\top-
\Sigma+W\preceq0$ for $K=L\Sigma^{-1}$; a Lyapunov-eigenvector argument makes every
feasible $A_{\mathrm{cl}}$ Schur. On this set the objective equals the stationary
cost $J(K)$, and completion of squares against the optimal $K^\star$ gives
\begin{equation}\label{eq:cos}
J(K)-\operatorname{tr}(PW)=\operatorname{tr}\!\big((K-K^\star)^\top
(R+B_\theta^\top PB_\theta)(K-K^\star)\,\Sigma_K\big),
\end{equation}
with $P$ the stabilizing control-DARE solution. Since $R+B_\theta^\top PB_\theta
\succ0$ and $\Sigma_K\succ0$, the right side vanishes iff $K=K^\star$, so the
minimizer $(\Sigma^\star,L^\star,Z_0^\star)$ is unique.

\emph{(A4) Smooth parameter entry.} The variables and cones do not depend on
$\theta$; $\theta$ enters only through $A_\theta,B_\theta$ (and, if co-designed,
$W,Q,R$), each $C^1$ in $\theta$, so $\nabla_\theta V=\langle\lambda^\star,
\partial_\theta(\text{data})\rangle$ once the multiplier $\lambda^\star$ is unique.

\explainnote{Verification of Eq.~\eqref{eq:cos}, spelled out step by step since it's stated without proof: write $A_K=A_\theta+B_\theta K$ and let $\Sigma_K\succ0$ solve the closed-loop Lyapunov equation $\Sigma_K=A_K\Sigma_K A_K^\top+W$, so $W=\Sigma_K-A_K\Sigma_K A_K^\top$. Using $J(K)=\operatorname{tr}(Q\Sigma_K)+\operatorname{tr}(RK\Sigma_K K^\top)=\operatorname{tr}[(Q+K^\top RK)\Sigma_K]$ (cyclic trace) and substituting $W$ into $\operatorname{tr}(PW)$ gives $\operatorname{tr}(PW)=\operatorname{tr}(P\Sigma_K)-\operatorname{tr}(A_K^\top PA_K\,\Sigma_K)$, hence
\[
J(K)-\operatorname{tr}(PW)=\operatorname{tr}\Big[\big(Q+K^\top RK+A_K^\top PA_K-P\big)\,\Sigma_K\Big].
\]
Everything now reduces to the pure matrix identity that the bracket equals $(K-K^\star)^\top(R+B_\theta^\top PB_\theta)(K-K^\star)$, with no trace or $\Sigma_K$ involved. Substitute the control DARE $P=Q+K^{\star\top}RK^\star+A_{K^\star}^\top PA_{K^\star}$ for $P$ inside the bracket and expand $A_K^\top PA_K$ and $A_{K^\star}^\top PA_{K^\star}$ using $A_K=A_\theta+B_\theta K$. Every cross term containing $A_\theta^\top PB_\theta$ can then be eliminated using the stationarity relation that defines $K^\star$ in the first place, $B_\theta^\top PA_\theta=-(R+B_\theta^\top PB_\theta)K^\star$ (equivalently $K^\star=-(R+B_\theta^\top PB_\theta)^{-1}B_\theta^\top PA_\theta$). What remains, with $\Phi:=R+B_\theta^\top PB_\theta$, is exactly
\[
K^\top\Phi K-K^{\star\top}\Phi K-K^\top\Phi K^\star+K^{\star\top}\Phi K^\star
\]
\[
\qquad=\;(K-K^\star)^\top\Phi(K-K^\star)
\]
by direct expansion, which is Eq.~\eqref{eq:cos}. So this is a standard LQR completion-of-squares identity (the same one behind the classical LQR "policy improvement" argument), not something novel or fragile to this paper -- it checks out.}

\explainnote{Roadmap for this (dense) proof of (A3), since it's the longest argument in the paper: (1) show both LMI multipliers must be rank-deficient and supported exactly on the kernel of the active constraint (Eq.~\eqref{eq:rank1}); (2) that forces every feasible dual point into a low-dimensional "sandwich" form (Eq.~\eqref{eq:lift}) parametrized by just two small matrices $\Theta,\Xi$; (3) plug that form into the KKT stationarity conditions (Eq.~\eqref{eq:kkt}) and the $B_\theta$-terms cancel, leaving a single Stein/Lyapunov equation (Eq.~\eqref{eq:stein}) whose unique solution is exactly the familiar Riccati $P$. That's the whole argument -- worth a one-sentence version of this up front for readers who just want the punchline ($S_{11}^\star=P$) without following every algebraic step.}
\emph{(A3) Unique dual, strict complementarity, and $S_{11}^\star=P$.} This is the
one delicate assumption; we prove it by \emph{constructing} the dual and showing it
is a single point. At the optimum both constraints are rank deficient. Using the
Lyapunov equality $\Sigma^\star-W=A_{\mathrm{cl}}\Sigma^\star A_{\mathrm{cl}}^\top$
and $A_\theta\Sigma^\star+B_\theta L^\star=A_{\mathrm{cl}}\Sigma^\star$,
\begin{equation}\label{eq:rank1}
M_1^\star=\begin{bmatrix}K^\star\\ I\end{bmatrix}\Sigma^\star
\begin{bmatrix}K^\star\\ I\end{bmatrix}^{\!\top}\!,\quad
M_2^\star=\begin{bmatrix}A_{\mathrm{cl}}\\ I\end{bmatrix}\Sigma^\star
\begin{bmatrix}A_{\mathrm{cl}}\\ I\end{bmatrix}^{\!\top}\!,
\end{equation}
each of rank $n$, with $\ker M_1^\star=\operatorname{range}[\,I;-K^{\star\top}]$ and
$\ker M_2^\star=\operatorname{range}[\,I;-A_{\mathrm{cl}}^\top]$. For PSD multipliers
$N\succeq0$ (of $M_1$) and $S\succeq0$ (of $M_2$), complementary slackness
$NM_1^\star=0$, $SM_2^\star=0$ force $\operatorname{range}N\subseteq\ker M_1^\star$
and $\operatorname{range}S\subseteq\ker M_2^\star$; a PSD matrix with range in
$\operatorname{range}(J)$ must factor as $J(\cdot)J^\top$, so
\begin{equation}\label{eq:lift}
\begin{aligned}
&N=\begin{bmatrix}I\\ -K^{\star\top}\end{bmatrix}\Theta
\begin{bmatrix}I\\ -K^{\star\top}\end{bmatrix}^{\!\top}\!,\quad
S=\begin{bmatrix}I\\ -A_{\mathrm{cl}}^\top\end{bmatrix}\Xi
\begin{bmatrix}I\\ -A_{\mathrm{cl}}^\top\end{bmatrix}^{\!\top}\!,\\
&\Theta\in\mathbb S^m_+,\ \Xi\in\mathbb S^n_+ .
\end{aligned}
\end{equation}
Thus the whole dual is carried by two cores $\Theta=N_{11}$, $\Xi=S_{11}$, with
$N_{12}=-\Theta K^\star$ and $S_{12}=-\Xi A_{\mathrm{cl}}$. The Lagrangian
$\mathcal L=\operatorname{tr}(Q\Sigma)+\operatorname{tr}(RZ_0)-\langle N,M_1\rangle
-\langle S,M_2\rangle$ has three stationarity conditions,
\begin{equation}\label{eq:kkt}
\begin{aligned}
&\partial_{Z_0}\!:\,\Theta=R,\quad
\partial_{L}\!:\,N_{12}=-B_\theta^\top S_{12},\\
&\partial_{\Sigma}\!:\,Q-N_{22}-S_{11}-S_{22}-(A_\theta^\top S_{12}+S_{12}^\top A_\theta)=0.
\end{aligned}
\end{equation}
The first fixes $\Theta=R$, hence $N_{12}=-RK^\star$ and $N_{22}=K^{\star\top}RK^\star$;
the second then reads
\begin{equation}\label{eq:star}
B_\theta^\top\Xi A_{\mathrm{cl}}=-RK^\star .
\end{equation}
Writing $A_\theta=A_{\mathrm{cl}}-B_\theta K^\star$ in $\partial_\Sigma$, the
relation \eqref{eq:star} cancels every $B_\theta$-cross term---explicitly
$A_\theta^\top\Xi A_{\mathrm{cl}}+A_{\mathrm{cl}}^\top\Xi A_\theta=
2A_{\mathrm{cl}}^\top\Xi A_{\mathrm{cl}}+2K^{\star\top}RK^\star$---and
$\partial_\Sigma$ collapses to the closed-loop Stein equation
\begin{equation}\label{eq:stein}
\Xi=Q+K^{\star\top}RK^\star+A_{\mathrm{cl}}^\top\Xi A_{\mathrm{cl}} .
\end{equation}
Since $A_{\mathrm{cl}}$ is Schur, \eqref{eq:stein} has a unique solution; the control
DARE is exactly $P=Q+K^{\star\top}RK^\star+A_{\mathrm{cl}}^\top PA_{\mathrm{cl}}$, so
$\Xi=P$. Hence the dual is the single point $S=[\,I;-A_{\mathrm{cl}}^\top]P[\,I;
-A_{\mathrm{cl}}^\top]^\top$ with $S_{11}^\star=P$ (and $N =[\,I; -K^{\star\top}]R[\,I;
-K^{\star\top}]^\top$). Strict complementarity follows by rank counting:
$\operatorname{rank}N=m$ ($\Theta=R\succ0$) and $\operatorname{rank}S=n$
($\Xi=P\succ0$ from $Q\succ0$) complete the ranks of $M_1^\star,M_2^\star$ (both $n$)
to full. This verifies (A1)--(A3); with (A4), Corollary~\ref{cor:lqg} follows.
\hfill$\square$

The identity $S_{11}^\star=P$ (dually $S_{11}^{\prime\star}=\Sigma^\varepsilon$) is
what makes the covariance-SDP gradient $\partial J_{\mathrm{cont}}^\star/\partial W=
-S_{11}^\star$ of \eqref{eq:gradQRcont} numerically identical to the Riccati gradient
$\partial\operatorname{tr}(PW)/\partial W=P$ used in
Section~\ref{sec:sim:complexity}.

%=======================================================================
\section{Existence of a subgradient under (A1)}\label{app:subgrad}
%=======================================================================

This appendix supplies the quantitative Lipschitz bound behind
Proposition~\ref{prop:subgrad}: when (A1) holds but (A2)/(A3) fail, the design still
admits a subgradient. (A1) alone makes
$V$ locally Lipschitz, so its Clarke subdifferential $\partial V(\theta)$ is nonempty,
and the returned dual $\partial_\theta L(z^\star,\lambda^\star,\theta)$ (with
$\lambda^\star=(\mu^\star,S^\star)$) is one of its elements.

Assume (A1) on a neighborhood of $\theta$: for each nearby $\theta'$ the inner
problem~\eqref{eq:generic} is convex with $C^1$ data, Slater holds (so strong duality
holds with attained primal and dual optima), and the optimal primal and dual sets stay
bounded. Two facts drive the argument.

\emph{(i) Dual feasibility does not depend on $\theta$.} A dual point is feasible iff it
lies in the fixed cone $\mathcal K^\ast$ (e.g.\ $S\succeq0$), which involves no $\theta$.
Hence any dual optimum is feasible at \emph{every} nearby design, and weak duality gives
\begin{equation}\label{eq:weak}
V(\theta'')\ \ge\ \min_z L(z,\lambda,\theta'')\quad\text{for every }\lambda\in\mathcal K^\ast .
\end{equation}

\emph{(ii) $L$ is Lipschitz in $\theta$ on the relevant set.} Since the data are $C^1$
and the optimal $z,\lambda$ stay bounded near $\theta$, there is a finite $M$ with
$\|\partial_\theta L(z,\lambda,\cdot)\|\le M$ over that bounded set.

\emph{Local Lipschitz bound.} Fix nearby $\theta,\theta'$ and put
$\lambda=\lambda^\star(\theta')$. By strong duality at $\theta'$,
$V(\theta')=\min_z L(z,\lambda,\theta')$; by \eqref{eq:weak} at $\theta$,
$V(\theta)\ge\min_z L(z,\lambda,\theta)$. Subtracting,
\begin{equation*}
V(\theta')-V(\theta)\ \le\ \min_z L(z,\lambda,\theta')-\min_z L(z,\lambda,\theta).
\end{equation*}
Let $z_\theta$ attain the second minimum. Then $\min_z L(z,\lambda,\theta')\le
L(z_\theta,\lambda,\theta')$ while $\min_z L(z,\lambda,\theta)=L(z_\theta,\lambda,\theta)$,
so the right-hand side is at most $L(z_\theta,\lambda,\theta')-L(z_\theta,\lambda,\theta)
\le M\|\theta'-\theta\|$ by (ii). Swapping $\theta\leftrightarrow\theta'$ gives the
reverse inequality, hence
\begin{equation}\label{eq:lip}
|V(\theta')-V(\theta)|\ \le\ M\,\|\theta'-\theta\| .
\end{equation}

\emph{Conclusion.} By \eqref{eq:lip} $V$ is locally Lipschitz, so (Rademacher) it is
differentiable almost everywhere and its Clarke subdifferential
$\partial V(\theta)=\operatorname{conv}\{\lim_k\nabla V(\theta_k):\theta_k\to\theta,\ V
\text{ differentiable at }\theta_k\}$ is nonempty, compact and convex. At each
differentiability point $\nabla V=\partial_\theta L$ by Theorem~\ref{prop:envelope}, so
passing to the limit $\partial_\theta L(z^\star,\lambda^\star,\theta)\in\partial
V(\theta)$: the dual the solver returns is a genuine subgradient. When (A2)--(A3) also
hold this set collapses to the single point $\{\nabla V(\theta)\}$. Where (A1) itself
fails---at a feasibility boundary where $V\to+\infty$, or where Slater or boundedness of
the dual set is lost---the bound \eqref{eq:lip} breaks and no finite subgradient need
exist. \hfill$\square$

\suggestnote{References.bib wasn't included in this review pass, so I couldn't verify the bibkeys resolve or check citation accuracy (author names, venues, years). The file header already flags this (keys from three separate sources: References.bib, ContEst\_related\_work.bib, and two suggested BibTeX entries at the end of this file) -- please do run bibtex end-to-end once more before submission and confirm no "undefined citation" warnings, especially for the two owned works (grigoriadis1998integrate, ramadan\_ekfkoopman) whose entries are marked "VERIFY" in the comments below.}
\bibliographystyle{plain}
\bibliography{References}

% ----------------------------------------------------------------------
% Suggested BibTeX for two works you already own (verify before use):
%
% @inproceedings{grigoriadis1998integrate,
%   author    = {Grigoriadis, Karolos M. and Zhu, Guoming and Skelton, Robert E.},
%   title     = {Optimal Redesign of Linear Systems / An Algorithm to Integrate
%                Structure and Control Design},
%   booktitle = {Proc.\ IEEE Conf.\ Decision and Control (CDC)},
%   year      = {1996}}   % VERIFY title, venue, year, pages
%
% @article{ramadan_ekfkoopman,
%   author  = {Ramadan, Mohammad S. and Anitescu, Mihai},
%   title   = {An eKF--Koopman Operator Approach for Tractable Stochastic
%              Optimal Control},
%   journal = {IEEE ...},
%   year    = {2024}}      % VERIFY title, journal, volume, pages, year
% ----------------------------------------------------------------------

\end{document}